\chapter{Pagination, Filtering, Sorting \& Field Selection}
\label{ch:pagination}

\section*{Executive Summary}
\addcontentsline{toc}{section}{Executive Summary}

Every list endpoint you will ever ship is a liability until you answer four
questions precisely: \emph{how does the client walk the collection, how does
the client narrow it, how does the client order it, and how much of each
resource does the client actually need?} These are the questions of
pagination, filtering, sorting, and field selection, and they are deceptively
deep. The naive answers --- \texttt{OFFSET 100000}, a \texttt{WHERE} clause
string-concatenated from query parameters, an \texttt{ORDER BY} on whatever
column name the client sent, and a response body containing every column the
ORM happened to load --- are the four horsemen of the reporting-API apocalypse
that this chapter's case study reconstructs in painful detail.

The chapter's spine is a single idea with three incarnations:
\textbf{position is fragile, value is stable}. Offset pagination addresses
rows by their \emph{position} in an ordering; positions shift under
concurrent writes, and reaching a deep position costs work proportional to
the position itself. Cursor and keyset pagination address rows by their
\emph{value} at the boundary of the previous page; values are stable under
inserts and deletes elsewhere in the ordering, and reaching a boundary value
costs an index lookup --- $O(\log n)$, not $O(n)$. We derive both costs
formally, then demonstrate the concurrency anomaly with a deterministic,
scripted writer/reader interleaving that reproduces duplicates and skips
under offset pagination while keyset pagination remains exact.

The second half of the chapter treats the query surface around pagination:
filtering and sorting. Here the governing principle is \textbf{never let the
client speak SQL}. Clients speak a small, documented, allow-listed grammar of
fields and operators; the server parses that grammar into a typed abstract
syntax tree and compiles the tree into parameterized SQL, rejecting unknown
fields, unknown operators, and type mismatches with a clean 400. Sorting gets
the same treatment, plus one iron rule: \emph{every ordering ends with a
unique tie-breaker}, because a non-deterministic \texttt{ORDER BY} silently
corrupts every pagination strategy built on top of it. We close with field
selection (\texttt{?fields=}) and expansion (\texttt{?expand=}), where the
naive implementation reintroduces the $N{+}1$ query problem at the API
boundary, and with the database craft that makes all of it fast: composite
indexes aligned with keyset predicates, covering indexes, and how to read an
\texttt{EXPLAIN} plan well enough to know whether you got the index right.

Everything is runnable offline. Five complete Python listings --- a seeded
dataset generator for the fictional Northwind Robotics parts catalog, a
FastAPI service implementing all three pagination strategies with correct
headers and HMAC-signed cursors, an allow-listed filter/sort compiler, a
benchmark harness, and the concurrency-anomaly demo --- form a laboratory you
can execute end to end on a Windows workstation with PowerShell in under a
minute.

\begin{keyconceptbox}
A page is not a \emph{place} in a list; it is a \emph{predicate} on values.
\texttt{OFFSET 150000} asks the database to compute and discard 150{,}000
rows and breaks when the list shifts. \texttt{WHERE (cost, id) > (4200,
918273)} asks the index to seek directly to a stable boundary and is immune
to writes anywhere else in the ordering. Design list endpoints around
boundary values, and design filter and sort surfaces around allow-lists
compiled to bound parameters --- never around client-supplied SQL fragments.
\end{keyconceptbox}

%% ============================================================================
\section{Learning Objectives \& Prerequisites}
\label{sec:ch8-objectives}

\subsection{Learning Objectives}

After completing this chapter and its lab, you will be able to:

\begin{enumerate}[label=\textbf{LO-\arabic*.}]
  \item Derive the asymptotic cost of offset pagination ($O(k)$ scanned rows
        at depth $k$) and of keyset pagination ($O(\log n)$ seek plus $O(p)$
        page fetch), and predict latency growth qualitatively from first
        principles.
  \item Formalize the offset-pagination concurrency anomaly --- duplicates
        and skips caused by inserts or deletes ahead of the reader's window
        --- and explain why keyset pagination is immune to it.
  \item Implement all three pagination strategies (offset, transparent
        keyset, opaque signed cursor) in FastAPI with correct HTTP semantics,
        including RFC~8288 \texttt{Link} headers with \texttt{rel} values
        \texttt{next}, \texttt{prev}, \texttt{first}, and \texttt{last}.
  \item Design cursor tokens that are opaque, tamper-evident (HMAC-SHA256),
        versioned, and expiring, and choose deliberately between
        \texttt{has\_more} and \texttt{total\_count} page metadata.
  \item Write a keyset predicate using tuple comparison ---
        \texttt{WHERE (sort\_col, id) > (?, ?)} --- and align a composite
        index with it so the predicate stays sargable.
  \item Build an allow-listed filter grammar (operators
        \texttt{eq/ne/gt/gte/lt/lte/in/contains}) that parses into a typed
        AST and compiles to parameterized SQL, rejecting anything outside
        the grammar.
  \item Enforce deterministic ordering: multi-key sorts, explicit NULL
        ordering, and a unique tie-breaker column appended to every sort.
  \item Implement sparse fieldsets and resource expansion without the
        $N{+}1$ query problem, using DataLoader-style batching at the API
        boundary.
  \item Read an \texttt{EXPLAIN QUERY PLAN} output and judge whether a
        composite or covering index actually serves a keyset query.
  \item Measure and report offset-versus-keyset latency growth with a
        reproducible benchmark protocol, and interpret the resulting curve.
\end{enumerate}

\subsection{Prerequisites}

This chapter assumes Chapters~0--6. Concretely:

\begin{itemize}
  \item \textbf{Chapter~0:} the latency-measurement discipline (percentiles,
        reproducible benchmark protocols, synthetic Northwind Robotics data).
  \item \textbf{Chapters~1--2:} HTTP request/response mechanics, header
        semantics, and consuming APIs from Python clients.
  \item \textbf{Chapter~3:} REST resource modeling --- collections,
        identifiers, and why representations are projections of state.
  \item \textbf{Chapter~4:} serialization contracts; the response shapes here
        are exactly the JSON contracts of that chapter under pagination.
  \item \textbf{Chapters~5--6:} client-side security hygiene and building
        FastAPI services with dependency-injected configuration, which this
        chapter's reference service extends \cite{fastapidocs,pydanticdocs}.
  \item \textbf{Working knowledge of SQL:} \texttt{SELECT}, \texttt{WHERE},
        \texttt{ORDER BY}, \texttt{LIMIT}/\texttt{OFFSET}, and what a B-tree
        index is. We teach \emph{index alignment} for pagination, but not
        index internals from zero; \emph{Designing Data-Intensive
        Applications} \cite{kleppmann2017ddia} is the recommended companion.
  \item \textbf{Environment:} Python 3.11+ on Windows with PowerShell, a
        virtual environment, and the dependencies from this chapter's
        appendix installed.
\end{itemize}

%% ============================================================================
\section{Deep Technical Mechanics \& Architectural Concepts}
\label{sec:ch8-mechanics}

\subsection{The Four Dimensions of Collection Design}
\label{subsec:ch8-four-dimensions}

A collection endpoint such as \texttt{GET /parts} is a function from client
intent to a projected, ordered, bounded subset of a relation. Decompose that
function into four orthogonal dimensions and every design decision in this
chapter has a home:

\begin{enumerate}
  \item \textbf{Pagination} --- which \emph{window} of the ordered
        collection? (offset, cursor, keyset)
  \item \textbf{Filtering} --- which \emph{subset} of the relation?
        (allow-listed predicates)
  \item \textbf{Sorting} --- which \emph{order}? (multi-key, deterministic)
  \item \textbf{Field selection} --- which \emph{projection} of each
        resource? (sparse fieldsets, expansion)
\end{enumerate}

The dimensions are orthogonal in design but \emph{not} in implementation:
filtering changes the relation being paginated, sorting defines what
``next'' means, and field selection changes what each row costs to
serialize. A change to any dimension between page one and page two of a
client's walk invalidates the walk --- a consistency requirement we will
formalize in Section~\ref{subsec:ch8-filter-consistency}.

\subsection{Offset Pagination and Its True Cost}
\label{subsec:ch8-offset-cost}

Offset pagination is the strategy everyone reaches for first:

\begin{minted}{sql}
SELECT id, sku, name, unit_cost_cents
FROM parts
ORDER BY unit_cost_cents, id
LIMIT 50 OFFSET 100000;
\end{minted}

The semantics are simple: sort the relation, discard the first $k$ rows,
return the next $p$. The problem is the word \emph{discard}. A B-tree index
on \texttt{(unit\_cost\_cents, id)} lets the engine find the \emph{start} of
the ordering in $O(\log n)$, but it provides no random access to the
$k$-th entry: index leaf entries do not carry their ordinal position. The
engine must therefore \textbf{visit} entries $1 \ldots k$ to count its way to
entry $k{+}1$, and in the general case --- when the query selects columns
not present in the index --- each visited entry costs a heap or clustered-index
lookup to fetch the row, only for the row to be thrown away.

\begin{definition}[Offset page-fetch cost]
Let $n$ be the number of rows matching the filter, $k$ the offset, $p$ the
page size, $c_s$ the per-row cost of scanning (and, where required, fetching)
a row that will be discarded, and $c_f$ the per-row cost of fetching and
returning a row that will be emitted. Then the expected cost of one offset
page fetch is
\[
T_{\mathrm{offset}}(k, p) \;=\; c_{\mathrm{seek}} \;+\; k \cdot c_s \;+\; p \cdot c_f
\;=\; O(\log n) + O(k) + O(p).
\]
With $k$ up to $n$, worst-case page cost is $O(n)$: \emph{the deepest page
costs as much as reading the entire table.}
\end{definition}

Two consequences follow. First, latency \emph{grows linearly with depth}, so
the pagination API has a built-in denial-of-wallet vector: a crawler (or a
buggy retry loop) walking to depth $n$ consumes
$\sum_{k} O(k) = O(n^2 / p)$ total row visits. Second, aggregate user
experience degrades exactly where your most valuable queries live ---
auditors, backfill jobs, and reporting exports are the clients that walk
deep.

\begin{examplebox}
Northwind Robotics ships 200{,}000 parts. A page of 50 at offset 0 costs one
index seek and 50 row fetches --- sub-millisecond on any hardware. The same
page at offset 150{,}000 must first scan past 150{,}000 index entries (and,
without a covering index, fetch 150{,}000 rows that are immediately
discarded). On the reference hardware used for this chapter's benchmark that
page costs ${\sim}6$~ms against ${\sim}0.07$~ms for the equivalent keyset
seek --- a ${\sim}90\times$ gap that widens linearly with depth.
Section~\ref{subsec:ch8-benchmark} gives the full measured table.
\end{examplebox}

\subsection{The Concurrency Anomaly, Formalized}
\label{subsec:ch8-anomaly-theory}

Cost is only half the indictment. Offset pagination is also \emph{incorrect}
under concurrent writes, in a way that no isolation level fixes, because the
anomaly spans multiple transactions (one per page request).

Model the collection as a sequence $S = \langle r_1, r_2, \ldots, r_n
\rangle$ ordered by the sort key. A reader requests pages of size $p$:
page $j$ is $S[jp \ldots jp{+}p{-}1]$. Between two page requests, a writer
commits. Define $\Delta_{\mathrm{ins}}(j)$ as the number of rows inserted
at positions $\le jp$ (ahead of or inside the window boundary) and
$\Delta_{\mathrm{del}}(j)$ the number of such rows deleted.

\begin{definition}[Offset window drift]
After writer activity ahead of the boundary, the next offset page returns
\[
S'\big[jp \ldots jp{+}p{-}1\big]
\;=\;
S\big[jp - \Delta \ldots jp - \Delta + p - 1\big],
\qquad \Delta = \Delta_{\mathrm{ins}}(j) - \Delta_{\mathrm{del}}(j).
\]
Every insertion ahead of the boundary ($\Delta > 0$) replays $\Delta$ rows
the reader already saw (\textbf{duplicates}); every deletion ($\Delta < 0$)
conceals $|\Delta|$ rows the reader never saw (\textbf{skips}).
\end{definition}

The reader did nothing wrong; each individual page was a correct answer to
the query as executed. The anomaly is a property of \emph{positional
addressing}: an offset is a coordinate in a coordinate system that the writer
keeps renumbering. Section~\ref{subsec:ch8-anomaly-demo} proves this with a
deterministic interleaving in which three concurrent front-of-list inserts
cause the offset walker to emit rows 5, 9, and 13 twice and never emit rows
18, 19, and 20 at all.

\begin{warningbox}
Read-committed isolation --- the default in PostgreSQL, and effectively what
SQLite gives each statement --- makes this \emph{worse} to reason about, not
better: every page request sees the latest committed state, so the window
boundary moves underneath the reader between every pair of requests.
Snapshot isolation across the whole walk would fix consistency but requires
holding a transaction (and its snapshot) open across an unbounded number of
HTTP requests --- an operational anti-pattern of its own. Keyset pagination
sidesteps the entire dilemma.
\end{warningbox}

\subsection{Cursor and Keyset (Seek) Pagination}
\label{subsec:ch8-keyset}

Keyset pagination replaces the positional coordinate with a \emph{value
coordinate}: the sort-key values of the last row of the previous page. The
next page is defined by a predicate, not a position:

\begin{minted}{sql}
SELECT id, sku, name, unit_cost_cents
FROM parts
WHERE (unit_cost_cents, id) > (4200, 918273)   -- the keyset
ORDER BY unit_cost_cents, id
LIMIT 50;
\end{minted}

\subsubsection{Tuple comparison semantics}

The predicate \texttt{(a, b) > (x, y)} is a \emph{row-value} (tuple)
comparison with lexicographic semantics:
\[
(a, b) > (x, y) \;\equiv\; a > x \;\lor\; \big(a = x \land b > y\big).
\]
This is exactly the condition ``strictly after the boundary row in the
$(a,b)$ ordering,'' including rows that tie on $a$ but follow on $b$. The
second disjunct is why the tie-breaker column is not optional: without
\texttt{id}, all rows sharing the boundary value of
\texttt{unit\_cost\_cents} would be either replayed or skipped. SQL
standardizes row-value comparison and both SQLite (3.15+) and PostgreSQL
implement it; where a dialect lacks it, the expanded disjunction above is
the portable equivalent --- write it exactly, because a hand-rolled
\texttt{a >= x AND b > y} is subtly wrong (it drops the $a = x \land b > y$
survivors of the tie on one side and replays $a > x$ ties on the other).

\subsubsection{Why the seek is O(log n)}

Given a composite index on \texttt{(unit\_cost\_cents, id)} whose column
order matches the sort order, the keyset predicate is \emph{sargable}: the
engine descends the B-tree directly to the first entry satisfying the tuple
predicate --- $O(\log n)$ comparisons --- then reads $p$ consecutive leaf
entries, already in order:

\begin{definition}[Keyset page-fetch cost]
\[
T_{\mathrm{keyset}}(p) \;=\; c_{\mathrm{seek}} \;+\; p \cdot c_f
\;=\; O(\log n) + O(p),
\]
\emph{independent of depth.} Page 1 and page 10{,}000 cost the same.
\end{definition}

The index-alignment requirement is precise: the index's leading columns must
equal the \texttt{ORDER BY} columns in the same order and direction (or all
reversed, which a B-tree can scan backward). An index on
\texttt{(id, unit\_cost\_cents)} is useless for this query; an index on
\texttt{(unit\_cost\_cents)} alone serves the seek on the first column but
forces an in-memory sort for the tie-break --- check with
\texttt{EXPLAIN} (Section~\ref{subsec:ch8-indexes}).

\subsubsection{Stable ordering is a correctness requirement}

Both strategies presuppose a \emph{total} order. If the sort key is not
unique, rows that tie can be emitted in different orders by different
executions of the same query (the optimizer is free to choose), and neither
offsets nor keysets mean anything stable. The rule is absolute:

\begin{keyconceptbox}
\textbf{Every paginated ordering must end in a unique, immutable tie-breaker
column} --- almost always the primary key. \texttt{ORDER BY unit\_cost\_cents}
is a bug waiting for a tie; \texttt{ORDER BY unit\_cost\_cents, id} is a
contract. The same rule makes keyset predicates total, makes offset walks at
least deterministic, and makes multi-column sorts (Section
\ref{subsec:ch8-sorting}) composable.
\end{keyconceptbox}

\subsubsection{Cursor pagination: keyset with an envelope}

\emph{Cursor pagination} is keyset pagination with the boundary values
packaged into an opaque token the server hands to the client:
\texttt{next\_cursor} in, \texttt{cursor} out. The underlying seek is
identical; what changes is the encapsulation. The server can now hide
internal column names, version the payload, sign it against tampering, and
expire it --- Section~\ref{subsec:ch8-cursor-design} develops the full
design space. Stripe, Slack, and most large public APIs converge on this
shape precisely because it decouples the client contract from the storage
schema.

\subsection{Measured Performance: Offset vs.\ Keyset}
\label{subsec:ch8-benchmark}

Theory predicts linear-versus-constant growth; measurement confirms it.
Listing~\ref{lst:ch8-benchmark} benchmarks single-page fetches (page size
50, median of 15 repetitions) at increasing depths over a 200{,}000-row
Northwind Robotics parts table in SQLite, backed by the composite index
\texttt{idx\_parts\_cost\_id}. Results from one run on the reference
workstation (Windows, Python 3.13, warm page cache):

\begin{table}[h]
  \centering
  \small
  \begin{tabularx}{\textwidth}{r r r r X}
    \toprule
    \textbf{Depth} & \textbf{Offset (ms)} & \textbf{Keyset (ms)} &
    \textbf{Speedup} & \textbf{Observation} \\
    \midrule
    0       & 0.060 & 0.062 & 1.0x  & Identical: offset 0 \emph{is} a seek. \\
    1{,}000   & 0.074 & 0.065 & 1.1x  & Scan cost still below timer noise. \\
    10{,}000  & 0.214 & 0.069 & 3.1x  & Linear regime clearly visible. \\
    50{,}000  & 0.937 & 0.079 & 11.9x & Offset now dominated by discards. \\
    100{,}000 & 1.853 & 0.076 & 24.4x & Keyset flat within noise. \\
    150{,}000 & 5.908 & 0.065 & 91.5x & Offset super-linear: cache pressure. \\
    \bottomrule
  \end{tabularx}
  \caption{Measured single-page latency by strategy and depth
           (median of 15; illustrative, machine-dependent; page size 50,
           $n = 200{,}000$ rows, SQLite with composite index).}
  \label{tab:ch8-benchmark}
\end{table}

\begin{figure}[h]
  \centering
  \begin{tikzpicture}
    \begin{axis}[
        width=0.86\textwidth,
        height=7.2cm,
        xlabel={Page depth (rows skipped)},
        ylabel={Median page-fetch latency (ms)},
        xmin=0, xmax=160000,
        ymin=0, ymax=6.5,
        xtick={0, 25000, 50000, 75000, 100000, 125000, 150000},
        scaled x ticks=false,
        x tick label style={/pgf/number format/fixed},
        legend pos=north west,
        legend style={font=\footnotesize},
        grid=major,
        grid style={gray!20},
      ]
      \addplot[color=packtorange, mark=*, thick] coordinates {
        (0, 0.060) (1000, 0.074) (10000, 0.214)
        (50000, 0.937) (100000, 1.853) (150000, 5.908)
      };
      \addlegendentry{OFFSET / LIMIT (linear growth)}
      \addplot[color=darkblue, mark=square*, thick] coordinates {
        (0, 0.062) (1000, 0.065) (10000, 0.069)
        (50000, 0.079) (100000, 0.076) (150000, 0.065)
      };
      \addlegendentry{Keyset seek (depth-independent)}
    \end{axis}
  \end{tikzpicture}
  \caption{Offset versus keyset page-fetch latency at increasing depth.
           Measured values from Listing~\ref{lst:ch8-benchmark}; illustrative
           and machine-dependent --- the \emph{shape}, not the absolute
           numbers, is the finding.}
  \label{fig:ch8-latency-curve}
\end{figure}

Read the curve as the two cost formulas made visible:
$T_{\mathrm{offset}}$ is the line $k \cdot c_s$, and $T_{\mathrm{keyset}}$
is the constant $c_{\mathrm{seek}} + p \cdot c_f$. The slight super-linearity
at depth 150{,}000 is cache behavior --- discarded rows still flow through
the buffer pool --- which makes the production case \emph{worse} than the
asymptotic one, since production offset queries rarely enjoy a covering
index.

\subsection{Choosing a Strategy: A Decision Tree}
\label{subsec:ch8-decision-tree}

\begin{figure}[h]
  \centering
  \begin{tikzpicture}[
      font=\footnotesize,
      decision/.style={diamond, draw=darkblue, thick, align=center,
                       inner sep=1.2pt, text width=2.9cm},
      outcome/.style={rectangle, rounded corners, draw=packtorange, thick,
                      align=center, text width=3.4cm, fill=packtorange!8},
      lbl/.style={note, midway},
      >=stealth,
      node distance=1.05cm and 1.5cm,
    ]
    \node[decision] (random) {Client needs random page access?\\(``jump to page 47'')};
    \node[decision, below left=of random] (shallow) {Depth bounded \& small?\\(admin UI, $<$ few 1000 rows)};
    \node[outcome, below=of shallow] (offok) {\textbf{Offset}\\ LIMIT/OFFSET + total\_count; simplest contract};
    \node[decision, below right=of random] (stable) {Writes concurrent with walks?\\(feeds, ledgers, exports)};
    \node[decision, below=of stable] (opaque) {Hide schema internals / expire sessions?};
    \node[outcome, below left=of opaque] (cursor) {\textbf{Opaque cursor}\\ signed, versioned keyset token};
    \node[outcome, below right=of opaque] (keyset) {\textbf{Transparent keyset}\\ after\_cost / after\_id params};
    \node[outcome, right=1.6cm of opaque] (offcap) {\textbf{Offset, capped}\\ enforce max depth, offer export job instead};

    \draw[->, thick] (random) -- node[lbl, above left] {no} (stable);
    \draw[->, thick] (random) -- node[lbl, above right] {yes} (shallow);
    \draw[->, thick] (shallow) -- node[lbl, left] {yes} (offok);
    \draw[->, thick] (shallow.east) to[out=0,in=110] node[lbl, right, pos=0.35] {no} (offcap.north);
    \draw[->, thick] (stable) -- node[lbl, left] {yes / any} (opaque);
    \draw[->, thick] (opaque) -- node[lbl, above left] {yes} (cursor);
    \draw[->, thick] (opaque) -- node[lbl, above right] {no} (keyset);
  \end{tikzpicture}
  \caption{Pagination-strategy decision tree. Random access is the only
           property offset provides that keyset cannot; it is also the
           property that costs $O(k)$ and drifts under writes.}
  \label{fig:ch8-decision-tree}
\end{figure}

The honest summary: offset is acceptable when the collection is small,
walks are shallow, and random access is a genuine product requirement ---
an internal admin grid is the canonical case. Everything user-facing,
machine-consumed, exported, or infinite-scrolled should be keyset, and the
cursor envelope should be opaque the moment the API is public
\cite{gehl2021apidesign}.

\subsection{Cursor Token Design}
\label{subsec:ch8-cursor-design}

\subsubsection{Opaque versus transparent}

A \emph{transparent} cursor exposes the boundary values as ordinary query
parameters (\texttt{?after\_cost=4200\&after\_id=918273}). Clients can
construct them, debug them, and bookmark them --- and they couple your public
contract to your physical schema forever. An \emph{opaque} cursor wraps the
same values in an encoded, signed token. The client treats it as an
unanalyzable string; the server retains freedom to rename columns, change
the sort, or rotate the encoding. For public APIs, opacity is the default
recommendation; for internal APIs where debuggability dominates, transparent
keysets are defensible.

\subsubsection{A reference token format}

The token used by the reference service (Listing~\ref{lst:ch8-service}) is
\[
\underbrace{\texttt{base64url}(\mathrm{JSON})}_{\text{payload}}
\;\texttt{.}\;
\underbrace{\operatorname{HMAC\text{-}SHA256}_{\text{32 hex chars}}(\text{payload})}_{\text{truncated tag}}
\]
with payload
\begin{minted}{json}
{"v":1,"exp":1735686000,"cost":4200,"id":918273}
\end{minted}

\begin{itemize}
  \item \textbf{Version (\texttt{v}).} Cursors outlive deploys. A version
        field lets the server reject (or migrate) tokens minted by a previous
        contract instead of misinterpreting them as corrupt.
  \item \textbf{Expiry (\texttt{exp}).} A cursor is a capability: it encodes
        a position in a potentially sensitive ordering. Expiry bounds replay
        windows, bounds how long a leaked token is useful, and gives you a
        principled answer to ``how long may a client pause mid-walk?''
  \item \textbf{Signature.} HMAC-SHA256 over the payload, truncated to 128
        bits, verified with a constant-time comparison
        (\texttt{hmac.compare\_digest}). Verification runs \emph{before} any
        payload field is trusted. Tamper handling is uniform: any malformed,
        forged, expired, or wrong-version token receives the same 400
        problem-details response, revealing nothing about which check failed
        beyond what the legitimate client needs to recover (restart the
        walk).
  \item \textbf{Key management.} The signing key is server-side
        configuration loaded from the environment (a secrets manager in
        production), never committed, never logged. The demo key in
        Listing~\ref{lst:ch8-service} is explicitly marked demo-only.
\end{itemize}

\begin{warningbox}
Unsigned or merely \emph{encoded} cursors (plain base64 JSON) are guessable
and forgeable: a client can mint a cursor pointing anywhere in your
ordering, bypassing any access control you keyed off earlier pages, and can
probe field values by binary-searching the boundary. Signing is not
hardening; it is the minimum bar. Similarly, never put user-input-derived
SQL fragments in the token --- the payload is decoded into a typed
structure and re-bound as parameters, exactly like any other external input.
\end{warningbox}

\begin{figure}[h]
  \centering
  \begin{tikzpicture}[
      font=\footnotesize,
      stage/.style={rectangle, draw=darkblue, thick, align=center,
                    minimum width=2.35cm, minimum height=1.05cm, fill=blue!4},
      decision/.style={diamond, draw=packtorange, thick, align=center,
                       inner sep=1pt, text width=1.9cm},
      >=stealth,
      node distance=0.55cm and 0.8cm,
    ]
    \node[stage] (fetch) {Fetch page $j$\\(\texttt{LIMIT p+1})};
    \node[stage, right=of fetch] (build) {Build payload\\$v$, $\exp$, boundary values};
    \node[stage, right=of build] (sign) {Sign\\HMAC-SHA256};
    \node[stage, right=of sign] (emit) {Emit token\\body + \texttt{Link} header};
    \node[stage, below=1.2cm of emit] (recv) {Client returns\\token verbatim};
    \node[decision, left=of recv] (verify) {Signature, version, expiry OK?};
    \node[stage, left=of verify] (seek) {Keyset seek from\\boundary values};
    \node[stage, fill=packtorange!12, draw=packtorange, below=0.7cm of verify]
      (reject) {400 problem-details\\``invalid cursor''};

    \draw[->, thick] (fetch) -- (build);
    \draw[->, thick] (build) -- (sign);
    \draw[->, thick] (sign) -- (emit);
    \draw[->, thick] (emit) -- (recv);
    \draw[->, thick] (recv) -- (verify);
    \draw[->, thick] (verify) -- node[note, above] {yes} (seek);
    \draw[->, thick] (verify) -- node[note, right] {no} (reject);
    \draw[->, thick, dashed] (seek.west) to[out=180,in=210]
      node[note, left, align=center] {next iteration\\(page $j{+}1$)} (fetch.south west);
  \end{tikzpicture}
  \caption{Cursor lifecycle: minting (top row) and redemption (bottom row).
           Verification precedes any use of payload fields; the one-extra-row
           fetch (\texttt{LIMIT p+1}) computes \texttt{has\_more} without a
           \texttt{COUNT(*)} query.}
  \label{fig:ch8-cursor-lifecycle}
\end{figure}

\subsection{Hypermedia Controls: RFC 8288 Link Headers}
\label{subsec:ch8-link-headers}

Page navigation is a hypermedia concern, and HTTP already has a standard
surface for it: the \texttt{Link} header \cite{rfc8288}. Each link is a
target URI plus a relation type:

\begin{minted}{http}
HTTP/1.1 200 OK
Link: </parts/offset?limit=50&offset=100>; rel="next",
      </parts/offset?limit=50&offset=0>; rel="first",
      </parts/offset?limit=50&offset=50>; rel="prev"
Content-Type: application/json
\end{minted}

The registered relations that matter here are \texttt{next}, \texttt{prev},
\texttt{first}, and \texttt{last}. Emitting them buys three things:
clients can be \emph{driven by the hypermedia} instead of constructing URLs
(a client that follows \texttt{rel="next"} never needs to know whether you
paginate by offset, keyset, or opaque cursor --- strategy migration becomes
a server-side decision); intermediaries and crawlers understand the
semantics; and the absence of \texttt{rel="next"} is an unambiguous,
parseable end-of-collection signal. Body-level fields (\texttt{next\_cursor})
and \texttt{Link} headers coexist peacefully --- the reference service emits
both --- but if you implement only one, prefer the header: it is the
standard, content-negotiation-friendly surface \cite{rfc9110}.

Two construction rules matter. First, links must be \emph{complete and
self-contained}: the \texttt{next} URI reproduces every parameter that
defined the current page (limit, filters, sort) so the walk is resumable
without client-side state. Second, \texttt{rel="last"} requires knowing the
total, which ties it to the cost discussion that follows --- many APIs
deliberately omit it.

\subsection{Page Metadata: \texttt{has\_more} versus \texttt{total\_count}}
\label{subsec:ch8-metadata}

Clients want to know two different things: ``is there another page?'' and
``how many rows are there in total?'' The first is cheap; the second is not.

\textbf{\texttt{has\_more}} costs one row: fetch \texttt{LIMIT p + 1}, emit
$p$ rows, and set \texttt{has\_more = true} if a $(p{+}1)$-th row exists.
The cursor lifecycle in Figure~\ref{fig:ch8-cursor-lifecycle} builds this in.

\textbf{\texttt{total\_count}} costs a \texttt{COUNT(*)} over the full
filtered relation. At small scale that is a fast index scan; at scale it is
an $O(n)$ operation \emph{per page request}, executed again and again by
every client on every page --- the database recomputes a number almost no
client acts upon, while holding buffer-pool pages that latency-sensitive
queries need. On a 200{,}000-row SQLite table the count alone costs
multiples of a page fetch; on a 200-million-row PostgreSQL table under
concurrent load it is the query your on-call learns to hate (Section
\ref{sec:ch8-casestudy}).

When product requirements genuinely need a total, choose deliberately:

\begin{itemize}
  \item \textbf{Opt-in counting} --- \texttt{?include\_total=true}, as in the
        reference service --- so casual clients never pay for it.
  \item \textbf{Count once, not per page} --- return the total only on page
        one of a walk, or in a separate \texttt{/parts/stats} resource.
  \item \textbf{Estimate} --- PostgreSQL's planner keeps reltuples estimates;
        \texttt{EXPLAIN} row estimates, or a periodically refreshed
        \texttt{COUNT} materialized into a stats table, give
        ``${\sim}200{,}000$ results'' good enough for UI badges.
  \item \textbf{Exact counts only over small filtered sets} --- enforce a
        threshold: compute the exact count only when a capped probe
        (\texttt{SELECT ... LIMIT 10001}) returns few enough rows.
\end{itemize}

\subsection{Filtering: A Safe Query-Parameter Grammar}
\label{subsec:ch8-filtering}

Filtering is where list endpoints most often become injection vehicles. The
failure pattern is always the same shape: the server treats query parameters
as \emph{SQL text} rather than as \emph{values}:

\begin{minted}{python}
# NEVER do this. One crafted ?min_cost=1 OR 1=1-- later...
sql = f"SELECT * FROM parts WHERE unit_cost_cents >= {min_cost}"
\end{minted}

The defense has three layers, all present in
Listing~\ref{lst:ch8-filter-compiler}:

\subsubsection{Layer 1: an explicit, allow-listed grammar}

Define the public surface as a mini-grammar, not as SQL:

\begin{minted}{text}
filter   := clause (";" clause)*
clause   := field ":" op ":" value
field    := "sku" | "name" | "category" | "cost" | "weight" | "active"
op       := "eq" | "ne" | "gt" | "gte" | "lt" | "lte" | "in" | "contains"
value    := scalar | scalar ("," scalar)*        ; comma list for "in"
sort     := ["-"] field ("," ["-"] field)*       ; "-" = descending
\end{minted}

Note what the grammar achieves by construction. The field set is a
\emph{public naming layer}: clients say \texttt{cost}; the column is
\texttt{unit\_cost\_cents}. Internal columns that must never be filtered or
leaked (password hashes, tenant keys, margin data) simply do not exist in
the grammar. The operator set is closed and semantics are documented:
\texttt{in} takes a comma list, \texttt{contains} is case-insensitive
substring on text fields only, ordering operators reject text fields, and
\texttt{contains} escapes \texttt{LIKE} wildcards in user input so
\texttt{\%} in a search string is data, not a pattern.

\subsubsection{Layer 2: parse to a typed AST, then validate}

Each clause becomes a typed node --- \texttt{Comparison},
\texttt{InList}, \texttt{Contains} --- with the value already coerced to the
field's declared Python type. Anything unparseable, unknown, or mistyped
raises \texttt{FilterError}, which the HTTP layer maps to a 400
problem-details response \cite{rfc9110}. Validation happens on the AST,
where context is rich, rather than on the SQL string, where context is gone.

\subsubsection{Layer 3: compile the AST to parameterized SQL}

Compilation resolves each public field name to its column via the allow-list
mapping --- \emph{never} via the input string --- emits a fixed skeleton of
SQL with \texttt{?} placeholders, and returns the values as a bound parameter
list. User input never touches SQL text; an \texttt{OR '1'='1} probe becomes
a string literal compared against a column and matches nothing. This is the
AST-based translation pattern: the grammar bounds what can be expressed, the
AST bounds what can be validated, and parameter binding bounds what can be
executed.

\begin{keyconceptbox}
Filter safety is an architecture, not a regex. \textbf{Allow-list the
grammar} (fields and operators), \textbf{parse to a typed AST} (validate
with full context, reject with 400s), and \textbf{compile with bound
parameters} (user values never become SQL text). Any one layer alone has a
bypass; the composition is robust.
\end{keyconceptbox}

\subsubsection{Filter consistency with pagination}
\label{subsec:ch8-filter-consistency}

A page is only meaningful relative to the filter and sort that defined the
walk. Two rules follow. First, the filter and sort parameters are part of
the \emph{cursor}: either they travel verbatim in every \texttt{Link} target
and the server rejects a mid-walk change, or their hash is embedded in the
signed cursor payload so a mismatched continuation fails fast instead of
silently returning pages from a different query. Second, keyset predicates
and filters compose by conjunction --- \texttt{WHERE <filters> AND
(unit\_cost\_cents, id) > (?, ?)} --- which keeps the composite index
usable when the filter's leading column aligns with the index (the
\texttt{idx\_parts\_cat\_cost\_id} index exists precisely so that
\texttt{category = ? AND (cost, id) > (?, ?)} seeks instead of scans).
Filtering by a column you did not index and then keyset-paginating the
result reintroduces the scan you worked to eliminate; check the plan.

\subsection{Sorting: Multi-Key Orderings Done Safely}
\label{subsec:ch8-sorting}

Sorting composes with everything above, so its rules are short and strict:

\begin{itemize}
  \item \textbf{Allow-list sortable fields and directions}, exactly as for
        filters. \texttt{?sort=-cost,sku} parses to
        \texttt{ORDER BY unit\_cost\_cents DESC, sku ASC}; an unknown field
        is a 400, and direction accepts only the two documented tokens.
        Direction must be a parsed token, not interpolated text ---
        \texttt{ORDER BY \{user\_input\}} is SQL injection wearing a trench
        coat.
  \item \textbf{Append the unique tie-breaker automatically.} The compiler
        in Listing~\ref{lst:ch8-filter-compiler} always ends the key list
        with \texttt{id ASC}. A client asking for \texttt{-cost} receives
        \texttt{ORDER BY unit\_cost\_cents DESC, id ASC}: deterministic,
        total, and keyset-compatible (the keyset tuple then spans
        \emph{all} sort keys: \texttt{(cost, sku, id) > (?, ?, ?)}).
  \item \textbf{Pin NULL ordering explicitly.} SQL leaves NULL placement in
        \texttt{ORDER BY} implementation-defined (SQLite: NULLs first on
        \texttt{ASC}; PostgreSQL: NULLs last). If any sortable column is
        nullable, state it: \texttt{ORDER BY discontinued\_at ASC NULLS
        LAST, id ASC} --- and make the keyset predicate NULL-aware
        (\texttt{(col IS NULL)} disjuncts), or exclude NULL rows from
        cursored endpoints. Ambiguous NULL handling is a silent
        page-boundary corruptor.
  \item \textbf{Descending keysets invert the predicate.} Walking
        \texttt{ORDER BY cost DESC, id DESC} pages forward with
        \texttt{(cost, id) < (?, ?)}. Mixed directions
        (\texttt{cost DESC, id ASC}) are legal but the tuple predicate no
        longer maps to one comparison --- expand to the disjunctive form.
        Prefer uniform direction for cursor endpoints.
  \item \textbf{Multi-column sorts need multi-column indexes.} The index
        column order must match the \texttt{ORDER BY} order left-to-right;
        an index on \texttt{(cost, sku, id)} cannot serve
        \texttt{ORDER BY sku, cost, id}.
\end{itemize}

\subsection{Field Selection and Expansion}
\label{subsec:ch8-field-selection}

\subsubsection{Sparse fieldsets}

\texttt{?fields=sku,name,cost} lets the client declare the projection it
actually uses. The win is threefold: smaller payloads (egress cost,
mobile bandwidth), less serialization CPU, and --- the quiet giant ---
\emph{covering-index eligibility}: \texttt{SELECT sku, unit\_cost\_cents}
with a restricted projection can be answered entirely from an index without
touching the table heap. The safety rules mirror filtering: allow-list
selectable fields, resolve through the public-name mapping, default to a
documented field set when the parameter is absent, and always include
\texttt{id} (clients need identity even when they forget to ask).

\subsubsection{Expansion and the N+1 boundary problem}

\texttt{?expand=category\_detail,supplier} asks the server to embed related
resources inline. The naive implementation issues one query for the page of
50 parts and then \emph{one supplier query per part} --- 51 round trips for
one HTTP request, the classic $N{+}1$ transplanted from ORM-land to the API
boundary. The fix is DataLoader-style batching applied to REST: collect the
distinct foreign keys of the page, issue \emph{one} \texttt{WHERE id IN
(...)} query per expanded relation, and stitch results in memory:

\begin{minted}{python}
def expand_suppliers(conn: sqlite3.Connection,
                     parts: list[dict]) -> dict[int, dict]:
    """One batched query for all suppliers referenced by this page."""
    keys = sorted({p["supplier_id"] for p in parts})
    placeholders = ", ".join("?" for _ in keys)
    rows = conn.execute(
        f"SELECT id, name, lead_time_days FROM suppliers "
        f"WHERE id IN ({placeholders})", keys,
    ).fetchall()
    return {r["id"]: dict(r) for r in rows}
\end{minted}

Batching converts $1 + N \cdot k$ queries into $1 + k$ for $k$ expanded
relations. Guard the expansion surface with the same allow-list discipline
(expandable relations are enumerated; each has a bounded batch size), cap
expansion depth (no \texttt{expand=a.b.c} transitive graphs in a REST
endpoint --- that way lies GraphQL, Chapter~12), and remember the trade-off:
expansion couples cacheability and payload size to the \emph{largest}
consumer's needs. Sparse fieldsets shrink responses; expansion grows them.
Offer both, and let \texttt{Vary} and your cache keys reflect the parameters
\cite{rfc9110}.

\subsection{Database Interaction: Indexes, EXPLAIN, Covering}
\label{subsec:ch8-indexes}

All of the above stands or falls on physical design:

\begin{itemize}
  \item \textbf{Composite index alignment.} For the ordering
        \texttt{(unit\_cost\_cents, id)}, the index
        \texttt{idx\_parts\_cost\_id ON parts(unit\_cost\_cents, id)} makes
        both the seek and the ordered page read index-native. The column
        order in the index must match the sort order; equality-prefixed
        filters extend the rule: \texttt{category = ? AND (cost, id) > (?,?)}
        wants \texttt{(category, unit\_cost\_cents, id)} --- equality columns
        first, range/sort columns after.
  \item \textbf{Reading EXPLAIN.} \texttt{EXPLAIN QUERY PLAN} on the keyset
        query reports \texttt{SEARCH parts USING INDEX idx\_parts\_cost\_id
        (unit\_cost\_cents>?...)} --- a \texttt{SEARCH}, i.e.\ a B-tree seek.
        The offset variant reports the same index but \emph{scans} through
        $k$ entries; at scale a planner may even choose
        \texttt{SCAN parts} plus an in-memory sort, which is the
        ``wrong index'' smell. Words to fear in any plan for a paginated
        endpoint: \texttt{SCAN} (full), \texttt{USE TEMP B-TREE FOR ORDER
        BY} (sort not served by an index), and a rows-examined estimate that
        grows with your offset.
  \item \textbf{Covering indexes.} If the page projection is
        \texttt{(id, sku, unit\_cost\_cents)}, an index on
        \texttt{(unit\_cost\_cents, id)} \emph{covering} \texttt{sku} (via
        \texttt{INCLUDE} in PostgreSQL, or a third index column in SQLite)
        answers the whole page from the index --- \texttt{USING COVERING
        INDEX} in the plan --- eliminating $p$ heap fetches per page and,
        for offset, $k$ wasted heap fetches per deep page.
  \item \textbf{Write amplification is the bill.} Every index is paid for on
        every insert and update. The two indexes in the reference schema are
        justified by two real query shapes; index-per-hope is how write
        latency dies.
\end{itemize}

%% ============================================================================
\section{Production-Ready Code Implementation}
\label{sec:ch8-code}

Five complete, runnable listings form the chapter's laboratory. All are
Python 3.11+, fully typed, PEP~8, and offline-first: the only datastore is
local SQLite (stdlib \texttt{sqlite3}); the only third-party packages are
FastAPI and its test client (which needs \texttt{httpx}). Everything runs on
Windows PowerShell in seconds. Setup:

\begin{minted}{powershell}
cd "C:\Training\Training Areas\Data jobs learning\APIs - Usage and development\code"
python -m venv .venv
.\.venv\Scripts\Activate.ps1
pip install fastapi pydantic httpx pytest
python gen_dataset.py --rows 200000 --db parts.db
\end{minted}

\subsection{Listing 1: Deterministic Dataset Generator}
\label{subsec:ch8-code-generator}

The generator (Listing~\ref{lst:ch8-generator}) builds the Northwind Robotics
parts catalog: 200{,}000 synthetic rows with a skewed cost distribution
(many cheap fasteners, few expensive actuators), deterministic per-row
content derived from a seeded PRNG, and the two composite indexes that
Section~\ref{subsec:ch8-indexes} justified: \texttt{idx\_parts\_cost\_id}
for the unfiltered keyset walk and \texttt{idx\_parts\_cat\_cost\_id} for
category-filtered walks. \texttt{ANALYZE} gives the query planner real
statistics. Generation is pure CPU and one bulk insert --- about 3.5 seconds
for 200{,}000 rows on the reference workstation; scale to a million rows by
passing \texttt{-{}-rows 1000000} if you want the benchmark's linear region
to hurt properly.

\begin{minted}[caption={Deterministic synthetic dataset generator (\texttt{gen\_dataset.py}).},label={lst:ch8-generator}]{python}
"""Deterministic synthetic dataset generator for Chapter 8.

Generates the fictional *Northwind Robotics* parts catalog into a local
SQLite database. Offline-first: no network, no third-party packages.

Usage (PowerShell):
    python gen_dataset.py --rows 200000 --db parts.db
"""

from __future__ import annotations

import argparse
import random
import sqlite3
import time
from pathlib import Path

SCHEMA = """
CREATE TABLE IF NOT EXISTS parts (
    id              INTEGER PRIMARY KEY,
    sku             TEXT NOT NULL UNIQUE,
    name            TEXT NOT NULL,
    category        TEXT NOT NULL,
    unit_cost_cents INTEGER NOT NULL,
    weight_grams    INTEGER NOT NULL,
    active          INTEGER NOT NULL,
    updated_at      TEXT NOT NULL
);
-- Composite index aligned with the keyset ordering (unit_cost_cents, id).
CREATE INDEX IF NOT EXISTS idx_parts_cost_id
    ON parts (unit_cost_cents, id);
-- Composite index aligned with (category, unit_cost_cents, id) keysets.
CREATE INDEX IF NOT EXISTS idx_parts_cat_cost_id
    ON parts (category, unit_cost_cents, id);
"""

CATEGORIES: list[str] = [
    "actuator", "sensor", "chassis", "drivetrain", "power",
    "controller", "fastener", "optic",
]

NAME_PARTS: list[str] = [
    "servo", "gyro", "lidar", "torque", "chassis", "rail", "bearing",
    "gearbox", "battery", "inverter", "encoder", "bracket", "harness",
]

NAME_SUFFIXES: list[str] = [
    "Alpha", "Beta", "Prime", "XL", "Micro", "HD", "Mk2", "Pro",
]


def generate_rows(count: int, seed: int) -> list[tuple]:
    """Build ``count`` deterministic synthetic part rows."""
    rng = random.Random(seed)
    rows: list[tuple] = []
    for i in range(1, count + 1):
        name = f"{rng.choice(NAME_PARTS)}-{rng.choice(NAME_SUFFIXES)}"
        category = rng.choice(CATEGORIES)
        # Skewed cost distribution: many cheap parts, few expensive ones.
        cost = int(rng.expovariate(1 / 4000.0)) + 50
        weight = rng.randint(2, 25_000)
        active = 1 if rng.random() < 0.93 else 0
        # Deterministic pseudo-timestamp derived from the row number only.
        minute = (i * 7) % (60 * 24 * 365)
        updated = f"2025-01-01T00:00:00Z+{minute:07d}m"
        rows.append((f"NWR-{i:07d}", name, category, cost, weight, active, updated))
    return rows


def build_database(db_path: Path, rows: int, seed: int) -> None:
    """Create (or replace) the SQLite database with ``rows`` parts."""
    if db_path.exists():
        db_path.unlink()
    conn = sqlite3.connect(db_path)
    try:
        conn.executescript(SCHEMA)
        started = time.perf_counter()
        data = generate_rows(rows, seed)
        with conn:
            conn.executemany(
                "INSERT INTO parts "
                "(sku, name, category, unit_cost_cents, weight_grams, "
                " active, updated_at) VALUES (?, ?, ?, ?, ?, ?, ?)",
                data,
            )
        conn.execute("ANALYZE")
        conn.commit()
        elapsed = time.perf_counter() - started
        count = conn.execute("SELECT COUNT(*) FROM parts").fetchone()[0]
        print(f"generated {count:,} rows into {db_path} in {elapsed:.2f}s")
    finally:
        conn.close()


def main() -> None:
    parser = argparse.ArgumentParser(description=__doc__)
    parser.add_argument("--rows", type=int, default=200_000)
    parser.add_argument("--db", type=Path, default=Path("parts.db"))
    parser.add_argument("--seed", type=int, default=42)
    args = parser.parse_args()
    if args.rows < 1:
        raise SystemExit("--rows must be >= 1")
    build_database(args.db, args.rows, args.seed)


if __name__ == "__main__":
    main()
\end{minted}

\subsection{Listing 2: Three Pagination Strategies in One FastAPI Service}
\label{subsec:ch8-code-service}

The service (Listing~\ref{lst:ch8-service}) implements the full design space
of Section~\ref{sec:ch8-mechanics} behind three endpoints that share one
ordering --- \texttt{(unit\_cost\_cents, id)} --- so their results can be
compared row-for-row:

\begin{itemize}
  \item \texttt{GET /parts/offset} --- \texttt{LIMIT}/\texttt{OFFSET} with
        RFC~8288 \texttt{Link} headers (\texttt{first}/\texttt{next}/
        \texttt{prev}, plus \texttt{last} only when the client opts into the
        count via \texttt{include\_total=true}).
  \item \texttt{GET /parts/keyset} --- transparent keyset pagination:
        \texttt{after\_cost}/\texttt{after\_id} parameters, the tuple
        predicate \texttt{WHERE (unit\_cost\_cents, id) > (?, ?)}, and the
        \texttt{LIMIT p+1} trick driving \texttt{has\_more} and the
        \texttt{rel="next"} link.
  \item \texttt{GET /parts/cursor} --- opaque cursors: versioned JSON
        payload with expiry, base64url-encoded, HMAC-SHA256-signed and
        verified with \texttt{hmac.compare\_digest} \emph{before} any payload
        field is read. Any malformed, forged, wrong-version, or expired
        token receives the same 400 problem-details body.
\end{itemize}

Limits are clamped server-side (\texttt{MAX\_LIMIT = 200}): an unbounded
\texttt{limit} is a resource-exhaustion invitation. The \texttt{\_self\_check}
entry point pages 250 rows through all three strategies offline via
\texttt{TestClient}, asserts byte-identical id sequences, and confirms a
forged cursor is rejected --- run it with
\texttt{python pagination\_service.py}.

\begin{minted}[caption={FastAPI service implementing offset, keyset, and HMAC-signed cursor pagination (\texttt{pagination\_service.py}).},label={lst:ch8-service}]{python}
"""Chapter 8 reference service: three pagination strategies over one table.

FastAPI application exposing the Northwind Robotics parts catalog with:

* ``GET /parts/offset``  -- classic LIMIT/OFFSET with RFC 8288 Link headers.
* ``GET /parts/keyset``  -- keyset (seek) pagination with transparent params.
* ``GET /parts/cursor``  -- opaque, HMAC-signed, versioned, expiring cursors.

All three order by ``(unit_cost_cents, id)`` which is backed by the
composite index ``idx_parts_cost_id`` created by ``gen_dataset.py``.

Run the self-check (offline, uses FastAPI TestClient):

    python gen_dataset.py --rows 50000 --db parts.db
    python pagination_service.py
"""

from __future__ import annotations

import base64
import hashlib
import hmac
import json
import os
import sqlite3
import time
from pathlib import Path
from typing import Any

from fastapi import FastAPI, Query, Request, Response
from fastapi.responses import JSONResponse

DB_PATH = Path(os.environ.get("PARTS_DB", "parts.db"))
# Demo-only key. In production load from a secrets manager, never hardcode.
CURSOR_SECRET = os.environ.get("CURSOR_SECRET", "northwind-robotics-demo-key")
CURSOR_VERSION = 1
CURSOR_TTL_SECONDS = 3600
MAX_LIMIT = 200
DEFAULT_LIMIT = 50

SELECT_COLS = "id, sku, name, category, unit_cost_cents, weight_grams, active"
ORDER_BY = "unit_cost_cents, id"

app = FastAPI(title="Northwind Robotics Parts API", version="1.0.0")


def get_db() -> sqlite3.Connection:
    conn = sqlite3.connect(DB_PATH)
    conn.row_factory = sqlite3.Row
    return conn


def row_to_part(row: sqlite3.Row) -> dict[str, Any]:
    return {
        "id": row["id"],
        "sku": row["sku"],
        "name": row["name"],
        "category": row["category"],
        "unit_cost_cents": row["unit_cost_cents"],
        "weight_grams": row["weight_grams"],
        "active": bool(row["active"]),
    }


def clamp_limit(limit: int) -> int:
    if limit < 1:
        return DEFAULT_LIMIT
    return min(limit, MAX_LIMIT)


def link_header(base_url: str, links: dict[str, str]) -> str:
    """Render an RFC 8288 Link header value from rel -> URL mappings."""
    return ", ".join(f'<{base_url}{target}>; rel="{rel}"' for rel, target in links.items())


# --------------------------------------------------------------------------
# Strategy 1: offset pagination
# --------------------------------------------------------------------------
@app.get("/parts/offset")
def list_parts_offset(
    request: Request,
    response: Response,
    limit: int = Query(DEFAULT_LIMIT, ge=1, le=MAX_LIMIT),
    offset: int = Query(0, ge=0),
    include_total: bool = Query(False),
) -> dict[str, Any]:
    limit = clamp_limit(limit)
    conn = get_db()
    try:
        rows = conn.execute(
            f"SELECT {SELECT_COLS} FROM parts "
            f"ORDER BY {ORDER_BY} LIMIT ? OFFSET ?",
            (limit, offset),
        ).fetchall()
        body: dict[str, Any] = {"items": [row_to_part(r) for r in rows]}
        if include_total:
            # Deliberately opt-in: COUNT(*) is a full index scan at scale.
            total = conn.execute("SELECT COUNT(*) FROM parts").fetchone()[0]
            body["total_count"] = total

        base = str(request.url.path)
        links: dict[str, str] = {"first": f"{base}?limit={limit}&offset=0"}
        if rows:
            links["next"] = f"{base}?limit={limit}&offset={offset + limit}"
        if offset > 0:
            prev_offset = max(0, offset - limit)
            links["prev"] = f"{base}?limit={limit}&offset={prev_offset}"
        if "total_count" in body:
            last_offset = max(0, ((body["total_count"] - 1) // limit) * limit)
            links["last"] = f"{base}?limit={limit}&offset={last_offset}"
        response.headers["Link"] = link_header(base, links)
        return body
    finally:
        conn.close()


# --------------------------------------------------------------------------
# Strategy 2: keyset (seek) pagination -- transparent parameters
# --------------------------------------------------------------------------
@app.get("/parts/keyset")
def list_parts_keyset(
    request: Request,
    response: Response,
    limit: int = Query(DEFAULT_LIMIT, ge=1, le=MAX_LIMIT),
    after_cost: int | None = Query(None, ge=0),
    after_id: int | None = Query(None, ge=1),
) -> dict[str, Any]:
    limit = clamp_limit(limit)
    conn = get_db()
    try:
        if after_cost is None or after_id is None:
            rows = conn.execute(
                f"SELECT {SELECT_COLS} FROM parts ORDER BY {ORDER_BY} LIMIT ?",
                (limit + 1,),  # fetch one extra row to compute has_more
            ).fetchall()
        else:
            # Tuple comparison keeps the predicate sargable on the
            # composite index (unit_cost_cents, id).
            rows = conn.execute(
                f"SELECT {SELECT_COLS} FROM parts "
                f"WHERE (unit_cost_cents, id) > (?, ?) "
                f"ORDER BY {ORDER_BY} LIMIT ?",
                (after_cost, after_id, limit + 1),
            ).fetchall()

        has_more = len(rows) > limit
        page = rows[:limit]
        body: dict[str, Any] = {
            "items": [row_to_part(r) for r in page],
            "has_more": has_more,
        }
        if has_more and page:
            last = page[-1]
            base = str(request.url.path)
            nxt = (
                f"{base}?limit={limit}"
                f"&after_cost={last['unit_cost_cents']}&after_id={last['id']}"
            )
            response.headers["Link"] = link_header(base, {"next": nxt})
        return body
    finally:
        conn.close()


# --------------------------------------------------------------------------
# Strategy 3: opaque signed cursors (HMAC-SHA256, versioned, expiring)
# --------------------------------------------------------------------------
def _b64url_encode(raw: bytes) -> str:
    return base64.urlsafe_b64encode(raw).rstrip(b"=").decode("ascii")


def _b64url_decode(text: str) -> bytes:
    padding = "=" * (-len(text) % 4)
    return base64.urlsafe_b64decode(text + padding)


def encode_cursor(last_cost: int, last_id: int) -> str:
    payload = {
        "v": CURSOR_VERSION,
        "exp": int(time.time()) + CURSOR_TTL_SECONDS,
        "cost": last_cost,
        "id": last_id,
    }
    raw = json.dumps(payload, separators=(",", ":"), sort_keys=True).encode()
    body = _b64url_encode(raw)
    tag = hmac.new(CURSOR_SECRET.encode(), body.encode(), hashlib.sha256).hexdigest()
    return f"{body}.{tag[:32]}"


class CursorError(ValueError):
    """Raised for any malformed, forged, or expired cursor."""


def decode_cursor(token: str) -> tuple[int, int]:
    try:
        body, tag = token.rsplit(".", 1)
    except ValueError as exc:
        raise CursorError("cursor is malformed") from exc
    expected = hmac.new(
        CURSOR_SECRET.encode(), body.encode(), hashlib.sha256
    ).hexdigest()[:32]
    if not hmac.compare_digest(expected, tag):
        raise CursorError("cursor signature mismatch")
    try:
        payload = json.loads(_b64url_decode(body))
    except (ValueError, UnicodeDecodeError) as exc:
        raise CursorError("cursor payload is not valid JSON") from exc
    if payload.get("v") != CURSOR_VERSION:
        raise CursorError("unsupported cursor version")
    if int(payload.get("exp", 0)) < int(time.time()):
        raise CursorError("cursor has expired")
    return int(payload["cost"]), int(payload["id"])


@app.get("/parts/cursor")
def list_parts_cursor(
    request: Request,
    response: Response,
    limit: int = Query(DEFAULT_LIMIT, ge=1, le=MAX_LIMIT),
    cursor: str | None = Query(None),
) -> dict[str, Any]:
    limit = clamp_limit(limit)
    position: tuple[int, int] | None = None
    if cursor is not None:
        try:
            position = decode_cursor(cursor)
        except CursorError as exc:
            return JSONResponse(
                status_code=400,
                content={
                    "type": "about:blank",
                    "title": "Invalid cursor",
                    "status": 400,
                    "detail": str(exc),
                },
            )

    conn = get_db()
    try:
        if position is None:
            rows = conn.execute(
                f"SELECT {SELECT_COLS} FROM parts ORDER BY {ORDER_BY} LIMIT ?",
                (limit + 1,),
            ).fetchall()
        else:
            rows = conn.execute(
                f"SELECT {SELECT_COLS} FROM parts "
                f"WHERE (unit_cost_cents, id) > (?, ?) "
                f"ORDER BY {ORDER_BY} LIMIT ?",
                (position[0], position[1], limit + 1),
            ).fetchall()
        has_more = len(rows) > limit
        page = rows[:limit]
        body: dict[str, Any] = {
            "items": [row_to_part(r) for r in page],
            "has_more": has_more,
        }
        if has_more and page:
            last = page[-1]
            token = encode_cursor(last["unit_cost_cents"], last["id"])
            body["next_cursor"] = token
            base = str(request.url.path)
            response.headers["Link"] = link_header(
                base, {"next": f"{base}?limit={limit}&cursor={token}"}
            )
        return body
    finally:
        conn.close()


def _self_check() -> None:
    """Offline verification: page through all three strategies and compare."""
    from fastapi.testclient import TestClient

    if not DB_PATH.exists():
        raise SystemExit(f"database {DB_PATH} missing; run gen_dataset.py first")
    client = TestClient(app)

    expected = get_db().execute(
        "SELECT id FROM parts ORDER BY unit_cost_cents, id LIMIT 250"
    ).fetchall()
    expected_ids = [r[0] for r in expected]

    # Offset: walk five pages of 50.
    seen: list[int] = []
    offset = 0
    for _ in range(5):
        res = client.get(f"/parts/offset?limit=50&offset={offset}")
        assert res.status_code == 200, res.text
        assert "next" in res.headers["Link"]
        seen.extend(item["id"] for item in res.json()["items"])
        offset += 50
    assert seen == expected_ids, "offset walk diverged"

    # Keyset: walk with transparent after_* parameters.
    seen = []
    params: dict[str, Any] = {"limit": 50}
    for _ in range(5):
        res = client.get("/parts/keyset", params=params)
        assert res.status_code == 200, res.text
        payload = res.json()
        seen.extend(item["id"] for item in payload["items"])
        last = payload["items"][-1]
        params = {
            "limit": 50,
            "after_cost": last["unit_cost_cents"],
            "after_id": last["id"],
        }
    assert seen == expected_ids, "keyset walk diverged"

    # Cursor: walk with opaque signed tokens; then tamper and expect 400.
    seen = []
    token: str | None = None
    for _ in range(5):
        params = {"limit": 50} if token is None else {"limit": 50, "cursor": token}
        res = client.get("/parts/cursor", params=params)
        assert res.status_code == 200, res.text
        payload = res.json()
        seen.extend(item["id"] for item in payload["items"])
        token = payload["next_cursor"]
    assert seen == expected_ids, "cursor walk diverged"

    forged = token[:-2] + ("00" if not token.endswith("00") else "11")
    res = client.get("/parts/cursor", params={"limit": 50, "cursor": forged})
    assert res.status_code == 400, res.text

    res = client.get("/parts/offset?limit=50&include_total=true")
    assert res.json()["total_count"] > 0
    assert "last" in res.headers["Link"]
    print("self-check OK: offset, keyset, and cursor strategies agree;" 
          " forged cursor rejected with 400")


if __name__ == "__main__":
    _self_check()
\end{minted}

\subsection{Listing 3: Allow-Listed Filter/Sort Compiler}
\label{subsec:ch8-code-compiler}

The compiler (Listing~\ref{lst:ch8-filter-compiler}) is the three-layer
defense of Section~\ref{subsec:ch8-filtering} as working code. The public
grammar (\texttt{field:op:value} clauses joined by \texttt{;}) parses into
frozen dataclass AST nodes; values are coerced to the field's declared type
at parse time; compilation resolves public names through the allow-list and
emits \texttt{?} placeholders plus a bound parameter list. The demo exercises
the happy path (a three-clause filter with a two-key descending sort), four
rejection classes (unknown field, unknown operator, type mismatch, malformed
clause), and an injection probe --- \texttt{name:eq:x' OR '1'='1} --- which
compiles to a bound string literal and correctly matches zero rows. Note the
public/private naming split: clients filter on \texttt{cost}, the column is
\texttt{unit\_cost\_cents}, and internal columns are simply absent from
\texttt{ALLOWED\_FIELDS}.

\begin{minted}[caption={Filter/sort compiler: allow-listed grammar, typed AST, parameterized SQL (\texttt{filter\_compiler.py}).},label={lst:ch8-filter-compiler}]{python}
"""Allow-listed filter/sort compiler: query params -> parameterized SQL.

Grammar (public, documented surface):

    ?filter=unit_cost_cents:gte:500;category:eq:sensor;name:contains:servo
    ?sort=-unit_cost_cents,sku        (leading '-' means descending)

The compiler never interpolates user input into SQL text. It builds a small
AST of typed nodes, then renders that AST to a WHERE clause plus a bound
parameter list. Unknown fields, unknown operators, and type mismatches are
rejected with a descriptive ``FilterError`` -- clients get a 400, never a 500.

Offline demo:  python filter_compiler.py
"""

from __future__ import annotations

import sqlite3
from dataclasses import dataclass
from pathlib import Path
from typing import Any, Literal

# The public allow-list: API name -> (column, python type).
# Internal column names are deliberately *not* the public names everywhere:
# "cost" is the public name for the internal column unit_cost_cents.
ALLOWED_FIELDS: dict[str, tuple[str, type]] = {
    "sku": ("sku", str),
    "name": ("name", str),
    "category": ("category", str),
    "cost": ("unit_cost_cents", int),
    "weight": ("weight_grams", int),
    "active": ("active", bool),
}

ALLOWED_OPERATORS = {"eq", "ne", "gt", "gte", "lt", "lte", "in", "contains"}
ORDERING_OPERATORS = {"gt", "gte", "lt", "lte"}


class FilterError(ValueError):
    """Raised for any filter/sort input outside the documented grammar."""


@dataclass(frozen=True)
class Comparison:
    field: str
    operator: Literal["eq", "ne", "gt", "gte", "lt", "lte"]
    value: Any


@dataclass(frozen=True)
class InList:
    field: str
    values: tuple[Any, ...]


@dataclass(frozen=True)
class Contains:
    field: str
    needle: str


@dataclass(frozen=True)
class SortKey:
    field: str
    descending: bool


FilterNode = Comparison | InList | Contains


def _coerce(field: str, raw: str) -> Any:
    column, kind = ALLOWED_FIELDS[field]
    del column  # coercion only needs the declared type
    try:
        if kind is bool:
            if raw.lower() not in {"true", "false"}:
                raise ValueError(raw)
            return 1 if raw.lower() == "true" else 0
        return kind(raw)
    except (TypeError, ValueError) as exc:
        raise FilterError(
            f"value {raw!r} is not a valid {kind.__name__} for field {field!r}"
        ) from exc


def parse_filter(expr: str) -> list[FilterNode]:
    """Parse the ``field:op:value[;field:op:value...]`` mini-grammar."""
    nodes: list[FilterNode] = []
    for clause in expr.split(";"):
        clause = clause.strip()
        if not clause:
            continue
        parts = clause.split(":", 2)
        if len(parts) != 3:
            raise FilterError(f"malformed clause {clause!r}; expected field:op:value")
        field, operator, raw = parts
        if field not in ALLOWED_FIELDS:
            raise FilterError(f"unknown field {field!r}; allowed: {sorted(ALLOWED_FIELDS)}")
        if operator not in ALLOWED_OPERATORS:
            raise FilterError(
                f"unknown operator {operator!r}; allowed: {sorted(ALLOWED_OPERATORS)}"
            )
        if operator == "in":
            values = tuple(_coerce(field, v) for v in raw.split(",") if v != "")
            if not values:
                raise FilterError("in operator requires at least one value")
            nodes.append(InList(field=field, values=values))
        elif operator == "contains":
            _, kind = ALLOWED_FIELDS[field]
            if kind is not str:
                raise FilterError(f"contains is only valid on text fields, not {field!r}")
            nodes.append(Contains(field=field, needle=raw))
        else:
            _, kind = ALLOWED_FIELDS[field]
            if kind is str and operator in ORDERING_OPERATORS:
                raise FilterError(f"operator {operator!r} is not valid on text field {field!r}")
            nodes.append(
                Comparison(field=field, operator=operator, value=_coerce(field, raw))
            )
    return nodes


def parse_sort(expr: str | None) -> list[SortKey]:
    """Parse ``-cost,sku`` into sort keys; always append the id tie-breaker."""
    keys: list[SortKey] = []
    if expr:
        for token in expr.split(","):
            token = token.strip()
            if not token:
                continue
            descending = token.startswith("-")
            field = token[1:] if descending else token
            if field not in ALLOWED_FIELDS:
                raise FilterError(f"cannot sort by unknown field {field!r}")
            keys.append(SortKey(field=field, descending=descending))
    # Deterministic tie-breaker: id is unique, so the total order is stable.
    keys.append(SortKey(field="id", descending=False))
    return keys


def compile_where(nodes: list[FilterNode]) -> tuple[str, list[Any]]:
    """Render filter nodes to a WHERE clause plus bound parameters."""
    fragments: list[str] = []
    params: list[Any] = []
    for node in nodes:
        column = ALLOWED_FIELDS[node.field][0]  # allow-list lookup, not input
        if isinstance(node, Comparison):
            symbol = {"eq": "=", "ne": "!=", "gt": ">", "gte": ">=", "lt": "<", "lte": "<="}[node.operator]
            fragments.append(f"{column} {symbol} ?")
            params.append(node.value)
        elif isinstance(node, InList):
            placeholders = ", ".join("?" for _ in node.values)
            fragments.append(f"{column} IN ({placeholders})")
            params.extend(node.values)
        elif isinstance(node, Contains):
            # Escape LIKE wildcards in user input; ESCAPE clause pins the marker.
            needle = node.needle.replace("\\", "\\\\").replace("%", "\\%").replace("_", "\\_")
            fragments.append(f"{column} LIKE ? ESCAPE '\\'")
            params.append(f"%{needle}%")
        else:  # pragma: no cover - exhaustive by construction
            raise FilterError(f"unsupported node {node!r}")
    clause = " AND ".join(fragments) if fragments else "1=1"
    return clause, params


def compile_query(
    filter_expr: str | None,
    sort_expr: str | None,
    limit: int,
) -> tuple[str, list[Any]]:
    """Compile public params into a complete, fully parameterized SELECT."""
    if limit < 1 or limit > 200:
        raise FilterError("limit must be between 1 and 200")
    nodes = parse_filter(filter_expr) if filter_expr else []
    where, params = compile_where(nodes)
    order = ", ".join(
        f"{ALLOWED_FIELDS.get(k.field, ('id',))[0]} {'DESC' if k.descending else 'ASC'}"
        for k in parse_sort(sort_expr)
    )
    sql = (
        "SELECT id, sku, name, category, unit_cost_cents, weight_grams, active "
        f"FROM parts WHERE {where} ORDER BY {order} LIMIT ?"
    )
    return sql, [*params, limit]


def _demo() -> None:
    db = Path("parts.db")
    if not db.exists():
        raise SystemExit("parts.db missing; run gen_dataset.py first")
    conn = sqlite3.connect(db)
    try:
        sql, params = compile_query(
            "cost:gte:500;category:in:sensor,actuator;name:contains:servo",
            "-cost,sku",
            limit=5,
        )
        print(f"SQL:    {sql}")
        print(f"params: {params}")
        for row in conn.execute(sql, params):
            print(row)

        for bad in (
            "password_hash:eq:x",   # unknown field (internal column leaked?)
            "cost:drop:1",          # unknown operator
            "cost:gte:abc",         # type mismatch
            "cost:eq:1;sort:evil",  # malformed clause
        ):
            try:
                compile_query(bad, None, 10)
            except FilterError as exc:
                print(f"rejected {bad!r}: {exc}")
            else:
                raise SystemExit(f"input {bad!r} should have been rejected")

        # Injection attempt is inert: it becomes a bound string parameter.
        sql, params = compile_query("name:eq:x' OR '1'='1", None, 5)
        rows = conn.execute(sql, params).fetchall()
        assert rows == [], "injection probe must match nothing"
        print("injection probe neutralized as bound parameter; demo OK")
    finally:
        conn.close()


if __name__ == "__main__":
    _demo()
\end{minted}

\subsection{Listing 4: Benchmark Harness}
\label{subsec:ch8-code-benchmark}

The harness (Listing~\ref{lst:ch8-benchmark}) measures what
Section~\ref{subsec:ch8-benchmark} plotted: median single-page latency for
offset and keyset fetches at depths from 0 to 150{,}000, page size 50, 15
repetitions, warm cache. Two details keep it honest. First, the keyset
boundary for each depth is discovered with a one-off positional query
(\texttt{boundary\_at}) so both strategies fetch the \emph{same} page.
Second, medians rather than means suppress scheduler noise. The printed
table is the source of Table~\ref{tab:ch8-benchmark} and
Figure~\ref{fig:ch8-latency-curve}; per Chapter~0's discipline, the numbers
are illustrative artifacts of one machine --- the deliverable is the curve
shape and your ability to reproduce it.

\begin{minted}[caption={Offset-versus-keyset latency benchmark harness (\texttt{benchmark.py}).},label={lst:ch8-benchmark}]{python}
"""Benchmark harness: offset vs keyset page-fetch latency at depth.

Measures the median wall-clock time to fetch one page of 50 rows at
increasing logical depths into the (unit_cost_cents, id) ordering.

Numbers are machine-dependent and therefore *illustrative*: the textbook
reports the shape of the curve, never a promised absolute latency.

Run:  python gen_dataset.py --rows 200000 --db parts.db
      python benchmark.py
"""

from __future__ import annotations

import sqlite3
import statistics
import time
from pathlib import Path

PAGE_SIZE = 50
REPEATS = 15
DEPTHS = [0, 1_000, 10_000, 50_000, 100_000, 150_000]

OFFSET_SQL = (
    "SELECT id, sku, name FROM parts "
    "ORDER BY unit_cost_cents, id LIMIT ? OFFSET ?"
)
KEYSET_SQL = (
    "SELECT id, sku, name FROM parts "
    "WHERE (unit_cost_cents, id) > (?, ?) "
    "ORDER BY unit_cost_cents, id LIMIT ?"
)


def time_query(conn: sqlite3.Connection, sql: str, params: tuple) -> float:
    """Return median seconds over REPEATS executions of one page fetch."""
    samples: list[float] = []
    for _ in range(REPEATS):
        start = time.perf_counter()
        rows = conn.execute(sql, params).fetchall()
        samples.append(time.perf_counter() - start)
    assert len(rows) == PAGE_SIZE
    return statistics.median(samples)


def boundary_at(conn: sqlite3.Connection, depth: int) -> tuple[int, int]:
    """Return the (cost, id) tuple of the row at logical position ``depth``."""
    row = conn.execute(
        "SELECT unit_cost_cents, id FROM parts "
        "ORDER BY unit_cost_cents, id LIMIT 1 OFFSET ?",
        (depth,),
    ).fetchone()
    if row is None:
        raise ValueError(f"depth {depth} beyond table size")
    return int(row[0]), int(row[1])


def main() -> None:
    db = Path("parts.db")
    if not db.exists():
        raise SystemExit("parts.db missing; run gen_dataset.py first")
    conn = sqlite3.connect(db)
    try:
        total = conn.execute("SELECT COUNT(*) FROM parts").fetchone()[0]
        print(f"rows={total:,}  page_size={PAGE_SIZE}  repeats={REPEATS}")
        print(f"{'depth':>8} | {'offset ms':>10} | {'keyset ms':>10} | speedup")
        print("-" * 48)
        for depth in DEPTHS:
            if depth + PAGE_SIZE >= total:
                continue
            cost, pid = boundary_at(conn, depth)
            t_offset = time_query(conn, OFFSET_SQL, (PAGE_SIZE, depth))
            t_keyset = time_query(conn, KEYSET_SQL, (cost, pid, PAGE_SIZE))
            ratio = t_offset / t_keyset if t_keyset else float("inf")
            print(
                f"{depth:>8,} | {t_offset * 1000:>10.3f} | "
                f"{t_keyset * 1000:>10.3f} | {ratio:>6.1f}x"
            )
    finally:
        conn.close()


if __name__ == "__main__":
    main()
\end{minted}

\subsection{Listing 5: Concurrency-Anomaly Demonstration}
\label{subsec:ch8-anomaly-demo}

Finally, the anomaly made tangible (Listing~\ref{lst:ch8-anomaly}). A
single-threaded, fully deterministic script interleaves a reader walking
four pages of five rows with a writer that inserts one \emph{cheap} part ---
landing at the front of the \texttt{(unit\_cost\_cents, id)} ordering ---
between every pair of page fetches. This is a schedule two real connections
can produce; scripting it in one process makes it reproducible offline. The
script runs both walkers against a pristine catalog (baseline) and against
the interleaved catalog, then diffs:

\begin{minted}{text}
offset: with concurrent inserts -> duplicates=[5, 9, 13], skipped=[18, 19, 20]
 keyset: with concurrent inserts -> duplicates=none, skipped=none
\end{minted}

Read the offset line against the drift formula of
Section~\ref{subsec:ch8-anomaly-theory}: each front insert
($\Delta_{\mathrm{ins}} = 1$ at every boundary) shifts the window back by
one, replaying one previously emitted row per page (the duplicates) and
pushing three tail rows past the final page (the skips). The keyset walk
anchors on values; front inserts are strictly \emph{before} the anchor and
therefore invisible to every subsequent predicate. Twenty rows, five-line
output, one permanent intuition.

\begin{minted}[caption={Deterministic writer/reader interleaving proving offset drift and keyset stability (\texttt{anomaly\_demo.py}).},label={lst:ch8-anomaly}]{python}
"""Concurrency-anomaly demo: offset pagination duplicates, keyset does not.

A deterministic, single-threaded script that *interleaves* a reader (walking
pages of a parts list) with a writer (inserting cheap parts at the front of
the (unit_cost_cents, id) ordering between page fetches). This is exactly the
schedule two concurrent connections can produce; scripting it in one process
keeps the demo reproducible and offline.

Run:  python anomaly_demo.py
"""

from __future__ import annotations

import random
import sqlite3
from dataclasses import dataclass, field

DB = ":memory:"
PAGE = 5
TOTAL_PAGES = 4


def seed_catalog(conn: sqlite3.Connection) -> None:
    conn.execute(
        "CREATE TABLE parts ("
        " id INTEGER PRIMARY KEY, sku TEXT, unit_cost_cents INTEGER)"
    )
    rng = random.Random(7)
    rows = [
        (f"NWR-{i:05d}", 1000 + i * 100) for i in range(1, 21)
    ]
    del rng  # rows are literal: fully deterministic
    conn.executemany(
        "INSERT INTO parts (sku, unit_cost_cents) VALUES (?, ?)", rows
    )
    conn.commit()


@dataclass
class WalkResult:
    seen_ids: list[int] = field(default_factory=list)

    @property
    def duplicates(self) -> list[int]:
        seen: set[int] = set()
        dupes: list[int] = []
        for pid in self.seen_ids:
            if pid in seen and pid not in dupes:
                dupes.append(pid)
            seen.add(pid)
        return dupes


def writer_insert_cheap_part(conn: sqlite3.Connection, n: int) -> None:
    """Insert one part cheaper than everything else (lands at the front)."""
    conn.execute(
        "INSERT INTO parts (sku, unit_cost_cents) VALUES (?, ?)",
        (f"NWR-NEW-{n:02d}", 10 + n),
    )
    conn.commit()


def walk_offset(conn: sqlite3.Connection, writer: bool) -> WalkResult:
    """Paginate with LIMIT/OFFSET while the writer shifts the window."""
    result = WalkResult()
    for page_no in range(TOTAL_PAGES):
        rows = conn.execute(
            "SELECT id FROM parts ORDER BY unit_cost_cents, id "
            "LIMIT ? OFFSET ?",
            (PAGE, page_no * PAGE),
        ).fetchall()
        result.seen_ids.extend(r[0] for r in rows)
        if writer:
            # Between requests, another client inserts a cheap part.
            writer_insert_cheap_part(conn, page_no)
    return result


def walk_keyset(conn: sqlite3.Connection, writer: bool) -> WalkResult:
    """Paginate with a (cost, id) seek predicate under the same writer."""
    result = WalkResult()
    position: tuple[int, int] | None = None
    for page_no in range(TOTAL_PAGES):
        if position is None:
            rows = conn.execute(
                "SELECT id, unit_cost_cents FROM parts "
                "ORDER BY unit_cost_cents, id LIMIT ?",
                (PAGE,),
            ).fetchall()
        else:
            rows = conn.execute(
                "SELECT id, unit_cost_cents FROM parts "
                "WHERE (unit_cost_cents, id) > (?, ?) "
                "ORDER BY unit_cost_cents, id LIMIT ?",
                (position[0], position[1], PAGE),
            ).fetchall()
        result.seen_ids.extend(r[0] for r in rows)
        if rows:
            position = (int(rows[-1][1]), int(rows[-1][0]))
        if writer:
            writer_insert_cheap_part(conn, 100 + page_no)
    return result


def main() -> None:
    for strategy, walker in (("offset", walk_offset), ("keyset", walk_keyset)):
        conn = sqlite3.connect(DB)
        seed_catalog(conn)
        baseline = walker(conn, writer=False)
        conn.close()

        conn = sqlite3.connect(DB)
        seed_catalog(conn)
        shifted = walker(conn, writer=True)
        dupes = shifted.duplicates
        # Rows of the original catalog the reader *skipped* entirely.
        skipped = sorted(set(baseline.seen_ids) - set(shifted.seen_ids))
        print(
            f"{strategy:>7}: with concurrent inserts -> "
            f"duplicates={dupes or 'none'}, skipped={skipped or 'none'}"
        )
        conn.close()

    print(
        "Conclusion: offset windows are positional, so front inserts replay "
        "rows;\nkeyset windows are value-anchored, so the walk stays stable."
    )


if __name__ == "__main__":
    main()
\end{minted}

\begin{examplebox}
\textbf{Verification record.} Every listing above was executed offline
before this chapter shipped: dataset generation (200{,}000 rows in
${\sim}3.5$~s), the service self-check (offset, keyset, and cursor walks
agree on 250 ids; forged cursor rejected with 400), the compiler demo (four
rejection classes plus a neutralized injection probe), the anomaly demo
(output above), and the benchmark (Table~\ref{tab:ch8-benchmark}). Reproduce
all of it with five PowerShell commands in Section~\ref{sec:ch8-lab}.
\end{examplebox}

%% ============================================================================
\section{Common Pitfalls \& Anti-Patterns}
\label{sec:ch8-pitfalls}

Each entry below names the failure, explains the mechanism, and states the
fix. Most production pagination incidents are one of these seven.

\subsection{Offset at Depth}
\textbf{Failure:} a list endpoint that is fast in testing and degrades
linearly in production as the table grows, because tests walk page 1 and
users (or crawlers) walk page 4{,}000. \textbf{Mechanism:}
$T_{\mathrm{offset}} = O(k)$ --- Section~\ref{subsec:ch8-offset-cost}; every
deep page re-scans and discards everything before it. \textbf{Fix:} keyset
or cursor pagination for any collection that can outgrow memory; if the
product genuinely needs random access, cap the depth
(\texttt{offset + limit <= 10000}) and route true exports to an async job.

\subsection{Unsigned, Guessable Cursors}
\textbf{Failure:} base64-encoded (but unsigned) cursors whose payload a
curious client edits --- replaying walks, jumping arbitrarily, probing the
ordering, or injecting unexpected types into payload fields.
\textbf{Mechanism:} encoding is not authentication; the server trusted
client-carried state. \textbf{Fix:} HMAC-sign the payload, verify with a
constant-time comparison \emph{before} parsing, version the payload, expire
it, and answer every failure mode with the same undifferentiated 400
(Section~\ref{subsec:ch8-cursor-design}).

\subsection{\texttt{total\_count} on Every Page}
\textbf{Failure:} each page response embeds an exact \texttt{COUNT(*)}, so
every request of every client's walk re-runs an $O(n)$ aggregate.
\textbf{Mechanism:} the count is computed even though the client needed it
once (or never); under concurrency it is also instantly stale, so the
precision was illusory anyway. \textbf{Fix:} default to \texttt{has\_more}
via \texttt{LIMIT p+1}; make totals opt-in, estimate from planner statistics
when an approximation suffices, and never compute \texttt{rel="last"} links
by default (Section~\ref{subsec:ch8-metadata}).

\subsection{Missing Tie-Breaker}
\textbf{Failure:} \texttt{ORDER BY unit\_cost\_cents} on a column full of
duplicates; pagination ``mostly works'' and occasionally replays or drops
rows near ties --- differently on each deploy, plan change, or vacuum.
\textbf{Mechanism:} ties make the order non-total, so neither offsets nor
keysets name a stable boundary. \textbf{Fix:} always append a unique,
immutable column: \texttt{ORDER BY unit\_cost\_cents, id}. Enforce it in the
sort compiler, not in documentation.

\subsection{Unbounded Page Sizes}
\textbf{Failure:} the server honors \texttt{?limit=1000000}; one client
allocates a megabyte-scale JSON array per request and starves the pool.
\textbf{Mechanism:} \texttt{limit} is a resource reservation request from an
untrusted peer. \textbf{Fix:} a hard server-side ceiling
(\texttt{MAX\_LIMIT = 200} in the reference service) with silent clamping or
a 400; document the ceiling in the contract.

\subsection{Filter Parameters Interpolated into SQL}
\textbf{Failure:} \texttt{f"WHERE col >= \{param\}"} or, worse,
\texttt{f"ORDER BY \{sort\_param\}"}. \textbf{Mechanism:} query parameters
are attacker-controlled text; interpolation makes them executable
(Section~\ref{subsec:ch8-filtering}). This remains the top API injection
vector in the wild. \textbf{Fix:} allow-listed grammar $\to$ typed AST
$\to$ bound parameters. There is no safe-regex shortcut, and ``internal
API'' is not a mitigation --- internal callers get compromised too.

\subsection{Exposing Internal Column Names}
\textbf{Failure:} filter/sort/field parameters accepted as raw column names,
leaking schema internals (\texttt{password\_hash}, \texttt{tenant\_key}) to
API consumers and coupling the public contract to physical refactors.
\textbf{Mechanism:} the public surface \emph{is} the schema. \textbf{Fix:} a
public naming layer resolved through an allow-list map
(\texttt{cost} $\to$ \texttt{unit\_cost\_cents}); unknown names get a 400
that lists the \emph{allowed} fields, never the schema.

\begin{warningbox}
A composite anti-pattern ties the list together: \emph{offset pagination
$+$ per-page \texttt{COUNT(*)} $+$ client-controlled sort on a non-unique
column}. Each is defensible alone in a demo; combined at scale they produce
an endpoint that is slow, incorrect under writes, and non-deterministic ---
simultaneously. When you inherit such an endpoint, migrate consumers to
cursors behind the \texttt{Link} header (clients following \texttt{rel="next"}
need not notice) rather than breaking the contract outright.
\end{warningbox}

%% ============================================================================
\section{Hands-On Production Lab \& Real-World Case Study}
\label{sec:ch8-lab}

\subsection{Lab: Three Strategies, One Dataset, Zero Excuses}

\textbf{Time:} 45--60 minutes. \textbf{Prerequisites:} the environment from
Section~\ref{sec:ch8-code}. All commands are PowerShell, run from the
chapter's \texttt{code} directory.

\subsubsection{Part A --- Build and verify the dataset (5 min)}

\begin{minted}{powershell}
python gen_dataset.py --rows 200000 --db parts.db
python -c "import sqlite3; c=sqlite3.connect('parts.db'); print(c.execute('SELECT COUNT(*), MIN(unit_cost_cents), MAX(unit_cost_cents) FROM parts').fetchone())"
\end{minted}

Expected: \texttt{generated 200,000 rows into parts.db} in a few seconds,
then a count of 200{,}000. Re-run with the same \texttt{-{}-seed}: the
database is byte-comparable deterministic in content (same rows, same
order).

\subsubsection{Part B --- Exercise all three strategies (15 min)}

\begin{minted}{powershell}
python pagination_service.py
\end{minted}

The self-check walks 250 ids through \texttt{/parts/offset},
\texttt{/parts/keyset}, and \texttt{/parts/cursor}, asserts the three id
sequences are identical, and forges a cursor to confirm the 400 path. Then
probe by hand with the server running
(\texttt{uvicorn pagination\_service:app --port 8000} in a second terminal):

\begin{minted}{powershell}
# First page, opaque cursor strategy; note the Link header and next_cursor.
(Invoke-WebRequest "http://localhost:8000/parts/cursor?limit=3").Headers["Link"]
(Invoke-WebRequest "http://localhost:8000/parts/cursor?limit=3").Content

# Tamper with the cursor: change the final character; expect HTTP 400.
Invoke-WebRequest "http://localhost:8000/parts/cursor?limit=3&cursor=<forged>"

# Keyset, transparently: choose your own boundary.
Invoke-RestMethod "http://localhost:8000/parts/keyset?limit=3&after_cost=5000&after_id=100000"
\end{minted}

Deliverable: in your notes, record the \texttt{Link} header of each
strategy and state which relations appear and why (\texttt{last} only under
\texttt{include\_total=true}; no \texttt{prev} for cursors --- why?).

\subsubsection{Part C --- Break filtering, then fail to (10 min)}

\begin{minted}{powershell}
python filter_compiler.py
\end{minted}

Confirm the four rejection classes print, then write three attack strings of
your own (an \texttt{OR 1=1} variant, a semicolon chain, a
\texttt{UNION SELECT}) and verify each either parses to an inert bound
parameter or is rejected at the grammar. Deliverable: the compiled SQL and
parameter list for one attack, annotated with the layer that neutralized it.

\subsubsection{Part D --- Reproduce the anomaly (10 min)}

\begin{minted}{powershell}
python anomaly_demo.py
\end{minted}

Expected output (deterministic):

\begin{minted}{text}
offset: with concurrent inserts -> duplicates=[5, 9, 13], skipped=[18, 19, 20]
 keyset: with concurrent inserts -> duplicates=none, skipped=none
\end{minted}

Deliverable: modify the writer to \emph{delete} rows ahead of the window
instead of inserting, predict the outcome using the drift formula
($\Delta < 0$ $\Rightarrow$ skips), and confirm. Then explain, in two
sentences, why no transaction isolation level on the reader's individual
page queries prevents this.

\subsubsection{Part E --- Measure the curve (10 min)}

\begin{minted}{powershell}
python benchmark.py
\end{minted}

Compare against Table~\ref{tab:ch8-benchmark}. Your absolute numbers will
differ (hardware, SQLite version, cache state); the \emph{shape} must not:
offset roughly linear in depth, keyset flat within noise. Deliverables:
(i) your table; (ii) re-run with \texttt{-{}-rows 1000000} and report where
offset crosses 50~ms; (iii) run \texttt{EXPLAIN QUERY PLAN} on both
\texttt{OFFSET\_SQL} and \texttt{KEYSET\_SQL} and quote the line that proves
the keyset is a \texttt{SEARCH} using \texttt{idx\_parts\_cost\_id}.

\subsection{Case Study: The 2 a.m. Reporting API}
\label{sec:ch8-casestudy}

\textit{Names and numbers are anonymized composites drawn from real
postmortems; the failure physics is exact.}

A mid-size robotics-parts distributor --- call it the production twin of
our fictional Northwind Robotics --- exposed an internal reporting API used
by the finance team's nightly reconciliation export. The endpoint was the
composite anti-pattern incarnate: \texttt{GET /api/ledger?offset=\&limit=1000},
sorted by \texttt{posted\_at} (non-unique, millisecond timestamps from a
batch importer), every response carrying \texttt{total\_count}, filters
accepted as \texttt{?where=...} fragments ``because only finance uses it.''

At 2 a.m.\ the export job began its walk over what had quietly become an
8-million-row ledger after an acquisition data load. The walk is the
$O(n^2/p)$ trap in production form: 8{,}000 page requests, each scanning an
average of 4 million rows to reach its window --- about $3.2 \times 10^{10}$
row visits for the night --- and each request \emph{also} running
\texttt{COUNT(*)} to fill \texttt{total\_count} and the \texttt{rel="last"}
link. Page fetches that cost 2~ms at offset 0 cost 40+ seconds by the
middle of the walk. The connection pool saturated; the health-check endpoint
--- served by the same pool --- began timing out; the orchestrator concluded
the service was dead and began killing and rescheduling pods, whose cold
caches made the next pages slower still. The morning dashboard showed a
cascading failure across services that shared the database; the root cause
was one \texttt{OFFSET} clause.

And the report was \emph{wrong anyway}. The batch importer was still
inserting back-dated rows ahead of the walker's window; the export contained
duplicated postings (finance caught ${\sim}11{,}000$ double-counted line
items in review) and a silent gap where a midnight delete had shifted a
window forward. Non-unique \texttt{posted\_at} ordering meant the duplicates
were not even deterministic.

The fix, deployed over the following sprint, maps one-to-one onto this
chapter: (1) keyset pagination over \texttt{(posted\_at, ledger\_id)} with a
composite index to match --- walk time fell from 9 hours to 11 minutes and
became flat with respect to depth; (2) \texttt{total\_count} removed from
pages, replaced by \texttt{has\_more} and a nightly-refreshed stats table
for the UI badge; (3) the \texttt{?where=} fragment replaced by an
allow-listed filter grammar compiled to bound parameters (an audit found
the fragment interface was reachable by any authenticated internal user);
(4) a deterministic tie-breaker added to every ordered endpoint. The
postmortem's action item worth stealing: \emph{``No list endpoint ships
without a named pagination strategy, a maximum depth or page size, and an
EXPLAIN plan in the PR description.''}

\begin{figure}[h]
  \centering
  \begin{tikzpicture}[
      font=\footnotesize,
      tl/.style={rectangle, draw=darkblue, thick, minimum height=0.62cm,
                 minimum width=1.55cm, align=center},
      wrt/.style={rectangle, draw=packtorange, thick, fill=packtorange!12,
                  minimum height=0.62cm, align=center},
      >=stealth,
    ]
    % --- offset timeline ---
    \node[note, anchor=east] at (-0.4, 0) {\textbf{offset walk}};
    \node[tl] (o1) at (0.8, 0) {page 1\\rows 1--5};
    \node[tl, right=0.25cm of o1] (o2) {page 2\\rows 6--10};
    \node[tl, right=0.25cm of o2] (o3) {page 3\\rows 11--15};
    \node[wrt, right=0.25cm of o3] (o4) {page 4\\rows 15--19};
    \node[note, right=0.3cm of o4, align=left] {row 15 \emph{again}:\\duplicate};
    \draw[->, thick] (o1) -- (o2); \draw[->, thick] (o2) -- (o3);
    \draw[->, thick] (o3) -- (o4);

    % writer events
    \node[wrt] (w1) at ($(o1.south east)!0.5!(o2.south west) + (0,-1.0)$) {insert\\cheap part};
    \node[wrt] (w2) at ($(o2.south east)!0.5!(o3.south west) + (0,-1.0)$) {insert\\cheap part};
    \node[wrt] (w3) at ($(o3.south east)!0.5!(o4.south west) + (0,-1.0)$) {insert\\cheap part};
    \node[note, anchor=east] at (-0.4, -1.3) {\textbf{writer}};
    \draw[->, packtorange, thick, dashed] (w1.north) -- (o2.south);
    \draw[->, packtorange, thick, dashed] (w2.north) -- (o3.south);
    \draw[->, packtorange, thick, dashed] (w3.north) -- (o4.south);
    \node[note] at ($(w2) + (0,-0.75)$) {each front insert shifts every later OFFSET window by $+1$};

    % --- keyset timeline ---
    \node[note, anchor=east] at (-0.4, -3.4) {\textbf{keyset walk}};
    \node[tl] (k1) at (0.8, -3.4) {page 1\\anchor $(c_5, i_5)$};
    \node[tl, right=0.25cm of k1] (k2) {page 2\\$(c, i) > $ anchor};
    \node[tl, right=0.25cm of k2] (k3) {page 3\\$(c, i) > $ anchor};
    \node[tl, right=0.25cm of k3] (k4) {page 4\\$(c, i) > $ anchor};
    \node[note, right=0.3cm of k4, align=left] {rows 16--20:\\exact, no replay};
    \draw[->, thick] (k1) -- (k2); \draw[->, thick] (k2) -- (k3);
    \draw[->, thick] (k3) -- (k4);
    \draw[->, packtorange, thick, dashed] (w1.south) to[out=-90,in=90] (k2.north);
    \node[note] at ($(k2.south)!0.5!(k3.south) + (0,-0.55)$)
      {inserts land \emph{before} the anchor value $\Rightarrow$ invisible to the predicate};
  \end{tikzpicture}
  \caption{Write-concurrency anomaly timeline. Above: offset windows are
           positional, so each front-of-ordering insert replays one row into
           the next page. Below: keyset windows are anchored on boundary
           \emph{values}; writes ahead of the anchor never enter the
           predicate. Compare Listing~\ref{lst:ch8-anomaly}'s measured
           output.}
  \label{fig:ch8-anomaly-timeline}
\end{figure}

%% ============================================================================
\section{Self-Assessment Exercises \& Deep-Dive Questions}
\label{sec:ch8-self-assessment}

Attempt each before consulting the references. Exercises marked [code]
extend the chapter's listings; [proof] items demand a derivation or a
counterexample; [design] items demand a defensible trade-off.

\begin{enumerate}[label=\textbf{E8.\arabic*.}]
  \item {[proof]} Starting from $T_{\mathrm{offset}}(k,p) = c_{\mathrm{seek}}
        + k c_s + p c_f$, show that a full walk of an $n$-row table in pages
        of $p$ costs $\Theta(n^2/p)$ row visits, and compute the constant
        factor. What does this imply for the \emph{worst} page size from the
        server's point of view?
  \item {[proof]} Give a precise writer schedule (inserts and deletes with
        positions) over a 20-row, page-size-5 walk such that the offset
        reader emits \emph{exactly} the multiset
        $\{r_1,\ldots,r_{19}\}$ with $r_{20}$ never emitted and $r_7$
        emitted twice.
  \item {[code]} Extend \texttt{anomaly\_demo.py} with a third walker that
        paginates by offset but re-reads the last row of the previous page
        to detect drift. Show it detects duplicates but cannot recover the
        skipped rows. Why is detection-without-recovery the best a purely
        positional scheme can do?
  \item {[code]} Add a \texttt{/parts/cursor} variant whose payload embeds a
        SHA-256 hash of the active filter and sort parameters, rejecting
        continuations whose parameters no longer match. Explain which attack
        or bug class this closes (Section~\ref{subsec:ch8-filter-consistency}).
  \item {[design]} Your product requires ``jump to page 47'' in an admin
        grid over 3 million rows. Design the endpoint: what do you cap, what
        do you count, what do you cache, and what do you tell the PM about
        the page-47 experience?
  \item {[proof]} Show that the tuple predicate \texttt{(a,b) > (x,y)} and
        its disjunctive expansion \texttt{a > x OR (a = x AND b > y)} are
        logically equivalent, and construct a 6-row table on which the
        hand-rolled predicate \texttt{a >= x AND b > y} returns a different
        result set. Name the two distinct error modes.
  \item {[code]} Add \texttt{NULLS LAST} ordering to a nullable
        \texttt{discontinued\_at} sort and write the NULL-aware keyset
        predicate for it. Verify against \texttt{filter\_compiler.py}'s
        demo database.
  \item {[design]} Design the expansion policy for \texttt{GET /orders
        ?expand=customer,lines.part}: batch sizes, depth cap, cache
        headers, and the point at which you tell the consumer to use the
        GraphQL gateway (Chapter~12) instead.
  \item {[code]} Using \texttt{EXPLAIN QUERY PLAN}, produce one index
        definition that turns the chapter's keyset query from
        \texttt{SEARCH ... USING INDEX} into \texttt{SCAN} plus a temporary
        B-tree, and explain in one paragraph what the planner is telling
        you.
  \item {[design]} A compliance rule requires an \emph{exact} total of
        matching records on an audited search endpoint. Reconcile this with
        Section~\ref{subsec:ch8-metadata}'s cost analysis; give at least
        two architectures and their failure modes.
\end{enumerate}

%% ============================================================================
\section{Interview Q\&A Vault}
\label{sec:ch8-interview-vault}

Ordered junior $\to$ staff. Answers are model responses at the depth a
strong candidate gives verbally.

\begin{enumerate}[label=\textbf{Q\arabic*.}, leftmargin=2.6em]
  % ---- junior ----
  \item \textbf{What is pagination and why can't we just return everything?}
        Pagination returns a collection in bounded pages. Unbounded
        responses blow up memory on both sides, serialize for seconds, blow
        through timeouts, and turn one curious client into a denial-of-service
        event. Bounding page size is a resource-governance decision as much
        as a UX one.
  \item \textbf{How does \texttt{LIMIT}/\texttt{OFFSET} work?}
        \texttt{ORDER BY} materializes the ordering, the engine skips
        \texttt{OFFSET} rows, then returns up to \texttt{LIMIT} rows. The
        skip is the problem: the engine must visit the skipped rows to count
        past them.
  \item \textbf{What is a cursor in API pagination?}
        An opaque token the server returns with a page; the client sends it
        back to get the next page. Internally it usually encodes the
        boundary row's sort-key values --- a bookmark by value, not by
        position.
  \item \textbf{What is the difference between \texttt{has\_more} and
        \texttt{total\_count}?}
        \texttt{has\_more} answers ``is there a next page'' and costs one
        extra row (\texttt{LIMIT p+1}). \texttt{total\_count} answers ``how
        many altogether'' and costs a full \texttt{COUNT(*)} over the
        filtered relation on every request that includes it.
  \item \textbf{What HTTP header is designed for page navigation links?}
        The \texttt{Link} header, RFC~8288 \cite{rfc8288}, with relation
        types \texttt{next}, \texttt{prev}, \texttt{first}, \texttt{last}.
        Clients can follow \texttt{rel="next"} without knowing the
        pagination strategy.
  \item \textbf{What does \texttt{?fields=id,name} do and why offer it?}
        Sparse field selection: the response includes only the requested
        attributes. It cuts payload size and serialization cost, and can
        make queries answerable from a covering index. Fields must come from
        an allow-list, not raw column names.
  \item \textbf{Why should every sort end with a unique column?}
        Because a non-unique sort key leaves ties in an unspecified order,
        and page boundaries through a tie are non-deterministic. Appending
        the primary key makes the order total, which both offset and keyset
        pagination require to be stable.
  \item \textbf{Name three safe filter operators you'd expose and one you
        would not.}
        \texttt{eq}, \texttt{gte}, \texttt{in} --- all parameterizable and
        index-friendly. I would not expose a raw \texttt{where} or
        free-text SQL fragment parameter: it cannot be made safe by
        sanitization, only by replacement with a grammar.
  % ---- mid ----
  \item \textbf{Why is offset pagination $O(n)$ at depth while keyset is
        $O(\log n)$?}
        A B-tree supports lookup by \emph{value}, not by \emph{position}.
        Reaching offset $k$ requires visiting $k$ index entries to count
        them --- $O(k)$, up to $O(n)$. A keyset predicate \texttt{(cost,id)
        > (?,?)} is a value lookup: $O(\log n)$ descent plus $O(p)$ for the
        page itself, independent of depth.
  \item \textbf{Explain the duplicates-and-skips anomaly of offset
        pagination.}
        Between page requests, inserts ahead of the reader's window shift
        every later position down by one, so the next \texttt{OFFSET}
        window replays rows already emitted (duplicates); deletes shift
        positions up, hiding unemitted rows (skips). Each page is
        individually correct; the composition is not, because the offset is
        a coordinate in a system the writer keeps renumbering.
  \item \textbf{Why doesn't snapshot isolation fix it?}
        It would --- if the whole walk ran in one transaction. But the walk
        spans unbounded HTTP requests; holding a snapshot open that long
        pins versions, bloats the database, and creates its own failure
        modes. Keyset anchoring fixes consistency without long-lived
        snapshots.
  \item \textbf{Write the keyset predicate for \texttt{ORDER BY cost, id}
        and expand the tuple comparison.}
        \texttt{WHERE (cost, id) > (?, ?)} --- equivalent to
        \texttt{cost > ? OR (cost = ? AND id > ?)}. The second disjunct is
        what keeps ties on \texttt{cost} correct, which is why the unique
        tie-breaker is part of the contract.
  \item \textbf{What index serves that query, and what must align?}
        A composite index on \texttt{(cost, id)} --- same column order and
        direction as the \texttt{ORDER BY}. An index on \texttt{(id, cost)}
        is useless for it; \texttt{(cost)} alone forces a sort of ties.
        Verify with \texttt{EXPLAIN}: you want \texttt{SEARCH ... USING
        INDEX}, not \texttt{SCAN} or a temp B-tree.
  \item \textbf{Design a tamper-proof cursor token.}
        JSON payload with version, expiry, and boundary values; base64url
        encode; append a truncated HMAC-SHA256 keyed with a server secret;
        verify with a constant-time comparison before parsing; on any
        failure return one undifferentiated 400. Payload fields are still
        re-bound as query parameters --- a valid signature does not make
        the contents trusted SQL.
  \item \textbf{Why version the cursor payload?}
        Tokens outlive deploys. A client mid-walk during a rollout can hold
        a token minted by the old build; a version field lets the new build
        reject it cleanly (client restarts the walk) instead of
        mis-parsing it as corrupt or, worse, valid-but-different.
  \item \textbf{When is offset pagination actually fine?}
        Small or slowly-changing collections, shallow walks, and a real
        random-access requirement --- an internal admin grid over ten
        thousand rows is the canonical case. Cap the maximum offset, and
        don't compute \texttt{COUNT(*)} per page.
  \item \textbf{How do you implement \texttt{rel="last"} and what does it
        cost?}
        You need the total: $last = \lfloor (n-1)/p \rfloor \cdot p$.
        That means \texttt{COUNT(*)} per request that emits it --- an
        $O(n)$ aggregate --- which is why many APIs omit \texttt{last} or
        compute totals from cached estimates.
  \item \textbf{How would you estimate \texttt{total\_count} cheaply?}
        Planner statistics (e.g.\ PostgreSQL \texttt{reltuples}), the row
        estimate from \texttt{EXPLAIN} of the same query, or a
        periodically refreshed stats table. All are approximate; present
        them as approximate (``${\sim}$200k results'') and never mix an
        estimate into a contract that promises exactness.
  \item \textbf{How do you keep client-supplied sort direction safe?}
        Direction is parsed as a token from a two-element set
        (\texttt{asc}/\texttt{desc} or a leading \texttt{-}), never
        interpolated. \texttt{ORDER BY \{user\_input\}} is SQL injection;
        the field allow-list protects the column, the direction token
        protects the keyword position, which bind parameters cannot reach.
  \item \textbf{How do filter parameters and cursors interact?}
        The filter and sort define the walk, so they're part of the cursor
        contract: either they must be resent verbatim on every page (and
        the server rejects mismatches) or their hash is embedded in the
        signed payload. A cursor silently continuing a \emph{different}
        query is a subtle correctness bug.
  \item \textbf{What is the N+1 problem at an API boundary and how do you
        fix it in REST?}
        \texttt{?expand=supplier} on a 50-item page becomes 1 query for the
        page plus 50 supplier queries. Fix with DataLoader-style batching:
        collect the page's foreign keys, run one \texttt{WHERE id IN (...)}
        per expanded relation, stitch in memory. $1{+}N$ queries become
        $1{+}k$ for $k$ relations.
  % ---- senior ----
  \item \textbf{How do keyset pagination and multi-column sort interact?}
        The keyset tuple spans \emph{all} sort keys in order ---
        \texttt{(a, b, id) > (?, ?, ?)} --- and the composite index must
        match column-for-column. Mixed sort directions break the single
        tuple comparison (you need the expanded disjunctive form with
        per-column operators), which is why cursor endpoints usually fix or
        constrain direction.
  \item \textbf{How do you paginate a feed ranked by score?}
        Score is the sort key: order by \texttt{(score DESC, id DESC)} and
        keyset on \texttt{(score, id) < (?, ?)}. The hard part is that
        scores \emph{change}: the walk then sees a time-varying ranking, so
        either snapshot the ranking (materialize a ranked list keyed by a
        snapshot id, cursor over positions \emph{within the immutable
        snapshot}) or accept drift and document it. Hot feeds (Reddit-style)
        snapshot; warm feeds accept drift.
  \item \textbf{Your keyset query suddenly regressed to seconds. How do you
        diagnose?}
        \texttt{EXPLAIN} first: look for \texttt{SCAN} replacing
        \texttt{SEARCH}, a temp B-tree for \texttt{ORDER BY}, or a
        rows-examined estimate that tracks the table size. Usual causes: a
        new filter column not prefixed in the index, a sort direction flip,
        stale planner statistics, or an OR-ed predicate defeating
        sargability. Then check whether the payload outgrew the covering
        index.
  \item \textbf{Transparent keyset vs.\ opaque cursor --- how do you
        choose?}
        Transparent keysets (\texttt{after\_cost=...}) are debuggable and
        bookmarkable but couple the public contract to the physical schema
        and let clients construct arbitrary boundaries. Opaque signed
        cursors decouple contract from schema, enable expiry and rotation,
        and prevent boundary forgery --- at the price of
        undebuggability-from-the-outside. Public API: opaque. Internal
        tools where engineers live in logs: transparent is defensible.
  \item \textbf{How do you paginate consistently across multiple database
        shards?}
        You don't get a global cheap ordering; you merge. Options: scatter
        keyset queries to all shards and $k$-way merge the pages (cost
        grows with shard count per page), pin a tenant to a shard so
        pagination is single-shard (the multi-tenant answer), or maintain a
        denormalized global ordering table. Offsets across shards compound
        the $O(k)$ scan $k$ ways --- keyset is essentially mandatory.
  \item \textbf{What are the security considerations of cursor design?}
        Cursors are capabilities: sign them (integrity), expire them
        (replay window), keep tenant/scope claims out unless enforced
        server-side per request (a valid cursor must not widen
        authorization), never log raw tokens, and treat payload fields as
        untrusted input even after signature verification. And the signing
        key rotates without invalidating the world only if you version keys.
  \item \textbf{How would you migrate a live public API from offset to
        cursors without breaking consumers?}
        Ship cursor support alongside offset on the same endpoint; emit
        \texttt{Link: rel="next"} in cursor form even for offset requests
        so hypermedia-driven clients migrate transparently; announce
        deprecation with a sunset header; add depth caps to offset first
        (they fail loudly for the crawlers you most want off it); measure
        strategy mix from logs; remove offset only at the next major
        version per the versioning policy (Chapter~9).
  \item \textbf{When does \texttt{COUNT(*)} actually become expensive, and
        what's the production alternative?}
        When the filtered relation stops fitting comfortably in cache or
        the count runs per page per client under concurrency --- at millions
        of rows it dominates endpoint latency. Alternatives: opt-in counts,
        counts on page one only, planner estimates, pre-aggregated stats
        tables refreshed out of band, or threshold counting (exact up to a
        cap, ``10000+'' beyond it, computed with a \texttt{LIMIT}ed probe).
  % ---- staff ----
  \item \textbf{Design pagination for a multi-tenant SaaS where tenant sizes
        span six orders of magnitude.}
        Per-tenant keyset over \texttt{(tenant\_id, sort\_key, id)} with a
        composite index led by \texttt{tenant\_id}: the equality prefix keeps
        every tenant's walk a local index range, so the 10-row tenant and
        the 100-million-row tenant share code without sharing performance
        failure modes. Depth caps and page-size ceilings apply per tenant;
        exports for whale tenants go to an async job with resumable cursors
        stored server-side.
  \item \textbf{How do you make a paginated export resumable across server
        restarts and client crashes?}
        Persist the cursor: the export job checkpoints the signed cursor
        (plus the filter/sort hash) after each page, so a crash resumes at
        the last boundary instead of page one. Keyset makes this sound ---
        the resume point is a value, immune to intervening writes; offset
        checkpoints resume into a shifted coordinate system and silently
        corrupt the export.
  \item \textbf{A client demands stable pagination over data being bulk
        backfilled. Options?}
        Four honest answers: (1) snapshot semantics --- materialize the
        query result or a snapshot id up front, paginate the frozen
        snapshot; (2) keyset with documented drift semantics (fine for
        appends at the tail, wrong for front inserts); (3) quiesce or
        fence the backfill behind a high-watermark column
        (\texttt{WHERE created\_at <= ?} baked into the walk) so the
        walked relation is effectively immutable; (4) change feed / CDC
        instead of pagination when the requirement is really ``see
        everything exactly once'' (Chapter~11).
  \item \textbf{Where do filtering capabilities belong architecturally ---
        gateway, service, or database?}
        Enforcement of the grammar and allow-list belongs at the service
        boundary (it is contract logic); compilation targets the database
        (it is execution logic); the gateway may enforce coarse limits
        (parameter count, page-size ceilings) but should not own filter
        semantics, or you cannot evolve the contract without redeploying
        infrastructure. The deeper the predicate is pushed into the query
        (sargable, index-aligned), the less data crosses every layer.
  \item \textbf{How do you test pagination correctness systematically?}
        Property tests, not examples: for random table contents, page
        sizes, and write schedules, assert the concatenation of pages equals
        the expected ordered multiset (keyset) or detect drift (offset).
        Add determinism tests (same walk twice $\Rightarrow$ same pages),
        boundary tests (empty table, single row, page size 1, exact
        multiples, ties spanning boundaries), security tests (forged,
        expired, and version-mismatched cursors), and a benchmark gate in
        CI so a missing index fails the build, not the 2 a.m.\ pager.
  \item \textbf{What consistency guarantees would you document for a public
        cursor-paginated endpoint, and why those?}
        Per-page serializability (each page is a consistent answer at some
        instant), monotonic value-progression (no page revisits a boundary
        it passed), and documented drift semantics for mutations of already-
        visited versus not-yet-visited regions. You explicitly do
        \emph{not} promise a stable total or a walk-level snapshot unless
        you sell snapshot cursors --- promising more than the mechanism
        delivers is how you end up in someone else's postmortem.
  \item \textbf{Is cursor pagination compatible with relational
        \texttt{JOIN}s and aggregation endpoints?}
        With care. The keyset predicate must reference columns of the
        driving relation and the index must cover the join's access path;
        joins multiply rows per key value, so the boundary tuple must be
        unique at the \emph{joined} row granularity or you split a group
        across pages. Aggregates generally should not paginate by keyset
        over grouped keys without a covering index on the grouping columns
        --- and pagination over a \texttt{GROUP BY} whose groups change
        under concurrent writes has the same drift physics as the base
        case, one level up.
  \item \textbf{What would you put in an organization-wide pagination
        standard?}
        One envelope shape (\texttt{items}, \texttt{has\_more},
        \texttt{next\_cursor}), one cursor format (signed, versioned,
        expiring), \texttt{Link} headers mandatory, \texttt{total\_count}
        opt-in only, \texttt{MAX\_LIMIT} set platform-wide, a unique
        tie-breaker on every ordered endpoint, an allow-listed filter
        grammar with a fixed operator set, an \texttt{EXPLAIN} plan
        attached to every list-endpoint PR, and a benchmark gate on
        depth-scaling. Standards exist so that the wrong thing is the hard
        thing.
  \item \textbf{Critique: ``We'll just paginate with
        \texttt{WHERE id > last\_id}''. When is that right, and when does it
        silently lie?}
        It is a keyset over the primary key --- optimal when the product
        ordering \emph{is} insertion order (ledgers, event logs, audit
        trails) because the clustered index serves it for free. It silently
        lies when consumers believe they're seeing a business ordering ---
        recency by \texttt{updated\_at}, price, relevance --- because id
        order is not any of those; it also breaks under id reassignment,
        merged imports with re-keyed ids, and any future requirement to
        re-sort. Right tool for append-only streams, wrong abstraction for
        ranked collections.
\end{enumerate}

%% ============================================================================
\section{External Bibliography \& Further Reading}
\label{sec:ch8-bibliography}

\begin{itemize}
  \item \textbf{RFC 8288, \emph{Web Linking}} \cite{rfc8288} --- the
        normative reference for \texttt{Link} headers and relation types
        (\texttt{next}/\texttt{prev}/\texttt{first}/\texttt{last}). Short,
        precise, and the reason hypermedia pagination is interoperable.
  \item \textbf{RFC 9110, \emph{HTTP Semantics}} \cite{rfc9110} --- status
        codes for the failure modes in this chapter (400 for grammar
        rejections), \texttt{Vary} for parameterized responses, and the
        caching semantics that interact with field selection.
  \item \textbf{Markus Winand, \emph{Use The Index, Luke}
        (\texttt{use-the-index-luke.com})} --- the canonical free reference
        on index anatomy; its ``Pagination Done the Right Way'' chapter is
        the keyset method's best-known exposition, with per-database
        dialect notes.
  \item \textbf{Martin Kleppmann, \emph{Designing Data-Intensive
        Applications}} \cite{kleppmann2017ddia} --- B-tree internals,
        snapshot isolation, and the read/write-path trade-offs behind every
        cost claim in this chapter.
  \item \textbf{JJ Geewax, \emph{API Design Patterns}}
        \cite{gehl2021apidesign} --- opinionated, battle-tested treatment of
        page tokens, filtering, and resource projection as repeatable
        design patterns.
  \item \textbf{PostgreSQL Documentation, \emph{LIMIT and OFFSET} and
        \emph{Indexes}} --- authoritative semantics for row skipping, row
        estimates used in counting strategies, and \texttt{NULLS
        FIRST/LAST} behavior.
  \item \textbf{SQLite Documentation, \emph{SELECT} and \emph{EXPLAIN QUERY
        PLAN}} --- row-value comparison support and the query-plan output
        vocabulary (\texttt{SEARCH} vs.\ \texttt{SCAN} vs.\ covering index)
        used throughout the lab.
  \item \textbf{JSON:API specification, \emph{Pagination and Sparse
        Fieldsets} profiles} --- a widely adopted convention family
        (\texttt{page[size]}, \texttt{page[after]}, \texttt{fields[type]})
        worth aligning with when you do not control both ends.
  \item \textbf{Stripe API documentation, \emph{Pagination}} --- the
        industry reference implementation of opaque cursor pagination
        (\texttt{starting\_after}/\texttt{ending\_before},
        \texttt{has\_more}); read it as a case study in minimal,
        stable contracts.
  \item \textbf{Slack API documentation, \emph{Pagination}} --- cursor
        pagination at scale (\texttt{next\_cursor} in
        \texttt{response\_metadata}), including the operational rationale
        for retiring offset-based methods.
  \item \textbf{FastAPI and Pydantic documentation}
        \cite{fastapidocs,pydanticdocs} --- query-parameter validation,
        response modeling, and the \texttt{TestClient} used by the
        reference service's offline self-check.
  \item \textbf{\emph{The Datadog/Shopify engineering blogs on cursor-based
        pagination and GraphQL Connections} (Relay Cursor Connections
        spec)} --- for the \texttt{edges}/\texttt{node}/\texttt{pageInfo}
        vocabulary you will meet again in Chapter~12.
\end{itemize}

%% ============================================================================
\section{Capstone Projects}
\label{sec:ch8-capstones}

Five projects, progressive. Each names its difficulty tier; acceptance
criteria are checkable facts, not vibes.

\subsection{Capstone 1: Cursor-Paginated Ledger API [junior]}

\textbf{Objective.} Re-implement the case study correctly: a
\texttt{/ledger} endpoint over a synthetic append-only postings table,
paginated by opaque signed cursor.

\textbf{Inputs/Datasets.} \texttt{gen\_dataset.py} extended with a
\texttt{ledger(posted\_at, ledger\_id, amount\_cents, account)} table,
seeded with 100{,}000 postings including deliberate timestamp ties.

\textbf{Steps.}
\begin{enumerate}
  \item Order by \texttt{(posted\_at, ledger\_id)}; build the composite
        index to match.
  \item Implement the keyset predicate with tuple comparison and the
        \texttt{LIMIT p+1} \texttt{has\_more} trick.
  \item Wrap the boundary in an HMAC-signed, versioned, expiring cursor;
        reject forged tokens with a uniform 400.
  \item Emit \texttt{Link: rel="next"} on every non-final page.
\end{enumerate}

\textbf{Acceptance Criteria.}
\begin{itemize}
  \item[$\square$] A scripted full walk emits every posting exactly once,
        verified by set comparison against the source table.
  \item[$\square$] \texttt{EXPLAIN QUERY PLAN} shows \texttt{SEARCH} using
        the composite index for every page query.
  \item[$\square$] Forged, expired, and version-mismatched cursors each
        return 400 with a problem-details body.
  \item[$\square$] Timestamp ties spanning a page boundary are handled
        without duplicates or skips (test included).
\end{itemize}

\textbf{Stretch Goals.} Add \texttt{rel="prev"} via a reversed-order
query; add per-account filtering with an index-led equality prefix.

\subsection{Capstone 2: Filter Compiler with Property Tests [mid]}

\textbf{Objective.} Harden Listing~\ref{lst:ch8-filter-compiler} into a
library with a documented grammar and adversarial test coverage.

\textbf{Inputs/Datasets.} The \texttt{parts} table; a corpus of 50 attack
strings (injection probes, wildcard floods, encoding tricks, type
confusions) you curate yourself.

\textbf{Steps.}
\begin{enumerate}
  \item Specify the grammar in the README-style module docstring (fields,
        operators, types, precedence-free conjunction).
  \item Extend the AST with parentheses-free \texttt{OR} groups
        (\texttt{|}-separated clauses) compiled with correct parentheses.
  \item Write \texttt{pytest} suites: unit tests per node type, golden
        SQL/params tests, and property tests asserting that (a) compiled
        SQL contains no substring of any raw input value, and (b) every
        attack string is either rejected or executes as inert data.
  \item Add a FastAPI adapter mapping \texttt{FilterError} to a 400
        problem-details response listing allowed fields and operators.
\end{enumerate}

\textbf{Acceptance Criteria.}
\begin{itemize}
  \item[$\square$] \texttt{pytest -q} green, including the property tests.
  \item[$\square$] All 50 attack strings rejected or provably inert
        (asserted by executing them against the dataset).
  \item[$\square$] No public filter/sort/field parameter name equals an
        internal column name it was not intended to expose.
  \item[$\square$] Compiled queries use bound parameters exclusively ---
        verified by inspection test on the compiler output.
\end{itemize}

\textbf{Stretch Goals.} Case-insensitive \texttt{contains} with correct
collation; filter-dependent index recommendation report driven by
\texttt{EXPLAIN}.

\subsection{Capstone 3: Pagination Benchmark Report [mid]}

\textbf{Objective.} Produce a reproducible, reviewer-grade benchmark
comparing offset, keyset, and cursor page fetches across table sizes and
page sizes.

\textbf{Inputs/Datasets.} \texttt{gen\_dataset.py} at 100k, 200k, and 1M
rows; \texttt{benchmark.py} extended with cursor-token overhead
measurement and p50/p95 columns.

\textbf{Steps.}
\begin{enumerate}
  \item Parameterize the harness by row count, page size, and depth list;
        record environment metadata (Python, SQLite version, CPU).
  \item Add percentile statistics (p50/p95) and a cold-cache variant using
        a fresh connection per repetition.
  \item Measure cursor mint/verify overhead per request and report it as a
        latency line item.
  \item Write the report: tables, one growth-curve figure per row count,
        and a stated verdict on where offset crosses your latency budget.
\end{enumerate}

\textbf{Acceptance Criteria.}
\begin{itemize}
  \item[$\square$] One command reproduces every table in the report
        (numbers flagged illustrative, shapes asserted).
  \item[$\square$] Offset growth is demonstrably linear (fit reported);
        keyset is flat within stated noise bounds.
  \item[$\square$] Cursor overhead is quantified and shown to be constant
        with depth.
  \item[$\square$] The report names the exact index used by each query and
        quotes the \texttt{EXPLAIN} evidence.
\end{itemize}

\textbf{Stretch Goals.} Covering-index variant showing the offset penalty
shrinking but not vanishing; PowerShell runner script for the whole matrix.

\subsection{Capstone 4: Infinite-Scroll Client over Signed Cursors [senior]}

\textbf{Objective.} Build a terminal-based infinite-scroll parts browser
that consumes \emph{only} \texttt{rel="next"} links, proving the client
never learns the pagination strategy.

\textbf{Inputs/Datasets.} The chapter's FastAPI service (cursor mode) over
the 200k-row dataset; \texttt{httpx} for the client.

\textbf{Steps.}
\begin{enumerate}
  \item Implement the client loop: follow \texttt{Link}
        \texttt{rel="next"} until absent; render pages with running totals
        and elapsed time per page.
  \item Add resume: persist the last cursor and the filter/sort state to
        disk; restart mid-walk and continue exactly.
  \item Add failure drills: expired cursor (sleep past TTL), forged cursor,
        server restart --- each handled with a documented recovery
        (restart walk vs.\ resume).
  \item Run the client against a live writer inserting rows mid-walk;
        record and explain the observable behavior.
\end{enumerate}

\textbf{Acceptance Criteria.}
\begin{itemize}
  \item[$\square$] Client source contains no offset arithmetic and no
        cursor parsing --- only header following.
  \item[$\square$] A kill-and-resume test completes the walk with zero
        duplicates (asserted client-side).
  \item[$\square$] Expired and forged cursors produce a controlled,
        user-visible recovery path, not a crash.
  \item[$\square$] Mid-walk inserts produce a written behavior note
        consistent with the chapter's drift model.
\end{itemize}

\textbf{Stretch Goals.} Conditional requests with \texttt{ETag} per page
(previews Chapter~10); strategy negotiation where the server silently
switches the client from offset links to cursor links.

\subsection{Capstone 5: Multi-Tenant Keyset Design [staff]}

\textbf{Objective.} Design and validate the pagination layer for a
multi-tenant SaaS whose tenants span six orders of magnitude in row count,
including isolation, index design, and org-wide standards.

\textbf{Inputs/Datasets.} A synthetic three-tenant dataset (1k, 200k, and
5M rows) in one \texttt{parts} table with \texttt{tenant\_id}; the
chapter's service as the starting point.

\textbf{Steps.}
\begin{enumerate}
  \item Add \texttt{tenant\_id} as the index-leading equality column:
        \texttt{(tenant\_id, unit\_cost\_cents, id)}; prove every tenant's
        walk stays a local \texttt{SEARCH}.
  \item Embed a tenant-scope claim \emph{outside} the cursor (server-side
        auth context) and prove a valid cursor from tenant A cannot page
        tenant B.
  \item Design per-tenant policy: page-size ceilings, depth caps, and an
        async export path with server-side resumable cursors for the whale
        tenant.
  \item Write the org standard as a one-page ADR: envelope shape, cursor
        format, error contract, and the EXPLAIN-in-the-PR gate.
  \item Load-test: concurrent walks from all three tenants; demonstrate
        the whale tenant's walk does not move the small tenants' p99.
\end{enumerate}

\textbf{Acceptance Criteria.}
\begin{itemize}
  \item[$\square$] Cross-tenant cursor replay is impossible, with a test
        proving it.
  \item[$\square$] All three tenants' page latency is depth-independent
        within stated noise, measured and tabulated.
  \item[$\square$] The whale tenant's export resumes correctly after a
        simulated crash (checkpoint test).
  \item[$\square$] The ADR names every trade-off explicitly rejected
        (offsets, per-page counts, transparent cursors) with reasons.
\end{itemize}

\textbf{Stretch Goals.} Row-level security in SQLite views or a second
connection role; snapshot-id cursors for compliance exports requiring a
frozen relation.

%% ============================================================================
\section{Python Dependencies Appendix}
\label{sec:ch8-dependencies}

The chapter's laboratory pins the following environment. Everything else is
the Python 3.11+ standard library --- notably \texttt{sqlite3},
\texttt{hmac}, \texttt{hashlib}, \texttt{base64}, and \texttt{random} ---
consistent with the curriculum's offline-first, stdlib-first rules.

\begin{minted}{text}
fastapi>=0.110
pydantic>=2.7
httpx>=0.27
pytest>=8
\end{minted}

Install on Windows PowerShell from the chapter's \texttt{code} directory:

\begin{minted}{powershell}
python -m venv .venv
.\.venv\Scripts\Activate.ps1
pip install -r requirements.txt
\end{minted}

\subsection{fastapi (\texttt{fastapi>=0.110})}

\textbf{Description.} The ASGI web framework used for the reference service
\cite{fastapidocs}: routing, request parsing, and automatic
query-parameter validation via type hints.

\textbf{Why included.} Chapters~6--7 standardized the book's API
implementations on FastAPI; this chapter's three pagination endpoints use
its typed \texttt{Query} parameters (with \texttt{ge}/\texttt{le} bounds
enforcing page-size ceilings) and its \texttt{TestClient} for the offline
self-check.

\textbf{Alternatives.} Flask (simpler, but validation is hand-rolled);
Django REST Framework (heavier, ORM-centric); a stdlib
\texttt{http.server} implementation (fine for demos, no validation
ecosystem).

\textbf{Installation.} \texttt{pip install "fastapi>=0.110"} --- pulls
\texttt{starlette} and \texttt{pydantic}; add \texttt{uvicorn} only if you
want to serve it live (\texttt{pip install uvicorn}).

\subsection{pydantic (\texttt{pydantic>=2.7})}

\textbf{Description.} Data-validation library \cite{pydanticdocs} powering
FastAPI's parameter coercion and the natural home for typed cursor-payload
models in larger codebases.

\textbf{Why included.} It is a FastAPI transitive dependency and the
contract-validation layer established in Chapter~4; this chapter's service
relies on its \texttt{Query(..., ge=1, le=200)} coercion to reject bad page
parameters before handler code runs.

\textbf{Alternatives.} \texttt{dataclasses} plus manual coercion (stdlib,
more code); \texttt{attrs} with validators; \texttt{marshmallow}.

\textbf{Installation.} \texttt{pip install "pydantic>=2.7"} --- installed
automatically with FastAPI; pin explicitly because v1/v2 semantics differ.

\subsection{httpx (\texttt{httpx>=0.27})}

\textbf{Description.} Modern HTTP client for Python, sync and async.

\textbf{Why included.} FastAPI's \texttt{TestClient} is built on
\texttt{httpx}; the offline self-check in
Listing~\ref{lst:ch8-service} and the infinite-scroll client in Capstone~4
both speak \texttt{httpx}. It also matches Chapter~2's client-side
toolchain.

\textbf{Alternatives.} \texttt{requests} (sync-only, ubiquitous);
\texttt{urllib} (stdlib, verbose); \texttt{aiohttp} (async, different API).

\textbf{Installation.} \texttt{pip install "httpx>=0.27"}.

\subsection{pytest (\texttt{pytest>=8})}

\textbf{Description.} The standard Python test runner and assertion
ecosystem.

\textbf{Why included.} Capstones~2--5 specify \texttt{pytest} suites
(property tests for the filter compiler, checkpoint and isolation tests for
the capstone services); it is the curriculum-wide test framework validated
by every track's \texttt{pytest -q} gate.

\textbf{Alternatives.} \texttt{unittest} (stdlib, no fixtures/plugins);
\texttt{nose2} (legacy); \texttt{hypothesis} as a \emph{complement} for
property-based testing if your track permits the extra dependency.

\textbf{Installation.} \texttt{pip install "pytest>=8"}.

%% ============================================================================
\section{Chapter Summary \& Key Takeaways}
\label{sec:ch8-summary}

\begin{itemize}
  \item \textbf{Position is fragile; value is stable.} Offset pagination
        costs $O(k)$ per page at depth $k$ and drifts (duplicates/skips)
        under concurrent writes; keyset pagination costs $O(\log n) + O(p)$
        at any depth and anchors on boundary values immune to writes
        elsewhere in the ordering.
  \item \textbf{The tuple predicate is the contract.} \texttt{WHERE
        (sort\_col, id) > (?, ?)} with a composite index in matching column
        order turns page $n$ into the same B-tree seek as page 1. Verify
        with \texttt{EXPLAIN}: \texttt{SEARCH ... USING INDEX}, never a
        temp B-tree sort.
  \item \textbf{Cursors are capabilities.} Opaque, HMAC-signed, versioned,
        expiring tokens, verified constant-time before parsing, rejected
        with a uniform 400. \texttt{Link} headers (RFC~8288) make
        navigation hypermedia-driven and strategy-migration possible.
  \item \textbf{Counting is not free.} Default to \texttt{has\_more}
        (\texttt{LIMIT p+1}); make \texttt{total\_count} opt-in, estimated,
        or out-of-band --- never a per-page $O(n)$ tax on every client.
  \item \textbf{Never let the client speak SQL.} Allow-listed grammar
        $\to$ typed AST $\to$ bound parameters, with public names decoupled
        from physical columns; direction tokens parsed, not interpolated.
  \item \textbf{Determinism is enforced, not hoped for.} Every ordering
        ends in a unique tie-breaker; NULL placement is stated explicitly;
        page sizes are capped server-side.
  \item \textbf{Projection and expansion are budget tools.} Sparse
        fieldsets shrink payloads and unlock covering indexes; expansion
        requires DataLoader-style batching or the $N{+}1$ problem moves to
        your API boundary.
\end{itemize}

%% ============================================================================
\section{Quick-Reference Card}
\label{sec:ch8-quickref}

\subsection*{Strategy Comparison Matrix}
\begin{table}[h]
  \centering
  \small
  \begin{tabularx}{\textwidth}{l c c c X}
    \toprule
    \textbf{Strategy} & \textbf{Cost at depth $k$} & \textbf{Write-stable?} &
    \textbf{Random access?} & \textbf{Use when} \\
    \midrule
    Offset (LIMIT/OFFSET) & $O(k)$ & No --- duplicates/skips & Yes &
      Small/shallow collections; admin grids; capped depth. \\
    Keyset (seek) & $O(\log n) + O(p)$ & Yes & No &
      Large collections; exports; feeds; internal APIs. \\
    Opaque cursor & $O(\log n) + O(p)$ & Yes & No &
      Public APIs: schema hiding, expiry, tamper evidence. \\
    \bottomrule
  \end{tabularx}
  \caption{Pagination strategies at a glance ($p$ = page size, $n$ = rows).}
\end{table}

\subsection*{Cursor Design Checklist}
\begin{enumerate}
  \item Payload: version field, expiry timestamp, boundary values ---
        nothing else.
  \item Integrity: HMAC-SHA256 (truncated), verified with
        \texttt{compare\_digest} \emph{before} parsing.
  \item Failures: malformed, forged, expired, wrong-version $\Rightarrow$
        one uniform 400 problem-details response.
  \item Key management: secret from environment/secrets manager; never
        logged; rotation plan via key version.
  \item Contract: \texttt{next\_cursor} in body \emph{and} \texttt{Link:
        rel="next"}; filter/sort state bound to the cursor (resent or
        hashed).
  \item Payload fields re-bound as SQL parameters --- signature $\neq$
        trust.
\end{enumerate}

\subsection*{Filter Grammar (Reference Implementation)}
\begin{minted}{text}
filter   := clause (";" clause)*
clause   := field ":" op ":" value
field    := "sku" | "name" | "category" | "cost" | "weight" | "active"
op       := "eq" | "ne" | "gt" | "gte" | "lt" | "lte" | "in" | "contains"
value    := scalar | scalar ("," scalar)*      # comma list for "in"
sort     := ["-"] field ("," ["-"] field)*     # "-" = descending
fields   := field ("," field)*                 # sparse fieldset
rules    : unknown field/op -> 400; values coerced to declared types;
           contains escapes LIKE wildcards; sort always appends ", id";
           SQL emitted with bound parameters only.
\end{minted}

\section{Standardized Evidence Addendum}
\subsection{Benchmark Snapshot Template}
\begin{table}[h]
  \centering
  \small
  \begin{tabularx}{\textwidth}{l c c c c}
    \toprule
    \textbf{Config} & \textbf{p50 latency} & \textbf{p99 latency} & \textbf{Error rate} & \textbf{Throughput} \\
    \midrule
    Baseline & -- & -- & -- & 1.00x \\
    Candidate A & -- & -- & -- & -- \\
    Candidate B & -- & -- & -- & -- \\
    \bottomrule
  \end{tabularx}
  \caption{Minimum benchmark table to keep per chapter.}
\end{table}

\subsection{Request Trace Template}
\begin{itemize}
  \item \textbf{Request:} one representative API call for this chapter's topic.
  \item \textbf{Before:} behavior (headers, payload, latency, failure mode) with the naive approach.
  \item \textbf{After:} behavior with the technique in this chapter.
  \item \textbf{Result:} what changed in correctness, resilience, latency, and operability.
\end{itemize}

\subsection{Cost and Latency Budget Snapshot}
\begin{table}[h]
  \centering
  \small
  \begin{tabularx}{\textwidth}{l c c}
    \toprule
    \textbf{Stage} & \textbf{Latency budget} & \textbf{Cost driver} \\
    \midrule
    TLS / connection setup & -- & handshake round trips, keep-alive hit rate \\
    Request validation & -- & schema complexity, CPU per request \\
    Business logic and I/O & -- & downstream fan-out, query cost \\
    Serialization and egress & -- & payload size, compression ratio \\
    \bottomrule
  \end{tabularx}
  \caption{Operational budget template for this chapter's technique.}
\end{table}

\section{Configuration Preset Template}
\begin{minted}{text}
# api_policy.yaml
policy_version: "v1.0.0"
component: "<chapter_topic>"
timeout_ms: 0
retry_budget: 0
fallback_strategy: "fail_fast"
rollback_trigger: "error_budget_burn_gt_threshold"
\end{minted}

\section{CI Quality Gate Specification}
\begin{itemize}
  \item contract tests green against the pinned OpenAPI/AsyncAPI/Protobuf spec,
  \item no breaking schema changes without a version bump,
  \item error responses conform to RFC 9457 problem-details shape,
  \item benchmark deltas within approved regression bounds (latency and throughput),
  \item policy version and rollback plan present in release metadata.
\end{itemize}

\section{Incident Runbook Template}
\begin{enumerate}
  \item Detect regression from SLO burn-rate alerts or consumer reports.
  \item Scope impacted endpoints, versions, and consumer segments.
  \item Contain impact (rollback, feature flag, rate-limit tightening, or traffic shift).
  \item Diagnose with traces, structured logs, and contract diffs.
  \item Recover with validated fix and verified rollout procedure.
  \item Prevent recurrence via new regression tests and CI gates.
\end{enumerate}

\section{Cross-Chapter Decision Card}
\begin{table}[h]
  \centering
  \small
  \begin{tabularx}{\textwidth}{l X X}
    \toprule
    \textbf{Dimension} & \textbf{Choose this approach when} & \textbf{Prefer an alternative when} \\
    \midrule
    Use case & -- & -- \\
    Latency budget & -- & -- \\
    Operational cost & -- & -- \\
    Team maturity & -- & -- \\
    \bottomrule
  \end{tabularx}
  \caption{Decision card to complete per chapter.}
\end{table}

\section{ROI Model and Build-vs-Buy}
\begin{itemize}
  \item \textbf{Build when:} the capability is differentiating, compliance forbids vendors, or scale makes vendor pricing irrational.
  \item \textbf{Buy when:} the capability is commodity (auth, gateways, observability), time-to-market dominates, or staffing is thin.
  \item \textbf{Hidden costs to model:} on-call load, upgrade cadence, expertise ramp, data egress fees.
\end{itemize}

\section{Disaster Recovery Notes (RPO/RTO)}
\begin{itemize}
  \item Define recovery point objective (RPO): maximum acceptable data loss window.
  \item Define recovery time objective (RTO): maximum acceptable restoration time.
  \item Test restores regularly; an untested backup is not a backup.
\end{itemize}

